% ===========================================================================
%  IntroAct-TS —— ICLR 2027 投稿正文
% ---------------------------------------------------------------------------
%  版本      : v4.5 review（2026-09-18，评审日志逐条落地版；占位符集记作 v46）
%  模板      : 用户提供的 ICLR 2027 style files（未作任何改动）
%  上一版    : introact_ts_iclr27_v44.tex
%
%  重要约定
%  ---------
%  1. 所有尚未在服务器上跑出的数字，一律写成占位符 \ph{KEY}，渲染为灰色
%     [KEY]。绝不允许用估计值、旧协议值或"看起来合理"的数填充。
%     占位符命名规范见 latex/README.md，完整清单见 docs/v46_placeholder_keys.csv
%     （1111 个唯一 KEY），执行侧规划见
%     docs/IntroActTS_20260918_v46_experiment_plan.md。
%  2. 只有 v4.3.1-r5 台账里已经独立审核过的历史数字，才允许以真实数值出现在
%     Appendix G，并且必须显式标注其属于旧协议、不得与 v4.5 主表比较。
%  3. 正文只保留 4.1-4.6 六个小节。其余全部进 Appendix。
%  4. 结论措辞受 docs/v431_r5_* 与规划 §41 的 claim 边界约束：在独立主实验完成
%     前不写 SOTA、不写 significantly outperforms、不写 safe/robust/generalizes。
%  5. 主表 roster 冻结为 KEEP / Best Fixed / R2-CART / T1 / TOI / TATO / SRDI
%     加 IntroAct-TS；ChannelTokenFormer 只作 external reference；Catalog Oracle
%     只作诊断。附录 C 的 30/50 与 Table 3 使用同一套七方法。
% ===========================================================================
\documentclass{article}

\usepackage{iclr2027_conference,times}
\input{math_commands.tex}

\usepackage{hyperref}
\usepackage{url}
\usepackage{graphicx}
\usepackage{booktabs}
\usepackage{amsmath}
\usepackage{amssymb}
\usepackage{multirow}
\usepackage{array}
\usepackage[table]{xcolor}

% ---------------------------------------------------------------- 版式与宏
\definecolor{phgray}{gray}{0.45}
\definecolor{bestgray}{gray}{0.94}

% 未运行结果的占位符。KEY 一律用下划线分层命名（如 MAIN_OVERALL_MASE），
% 下划线是数学下标字符，因此必须 \detokenize 成普通字符后才能进 \textsf。
% 填充时只替换 KEY，不要改 \ph 本身。
\newcommand{\ph}[1]{\textcolor{phgray}{\textsf{[\detokenize{#1}]}}}
% 表格中强调：第一名加粗，第二名下划线，Oracle 用斜体灰字。
\newcommand{\best}[1]{\textbf{#1}}
\newcommand{\second}[1]{\underline{#1}}
\newcommand{\oracle}[1]{\textit{\textcolor{phgray}{#1}}}
\newcommand{\introact}{\textsc{IntroAct-TS}}
\newcommand{\tsfm}{TSFM}
\newcommand{\calA}{\mathcal{A}}
\newcommand{\calB}{\mathcal{B}}
\newcommand{\calN}{\mathcal{N}}
\newcommand{\loss}{\ell}

\newcommand{\fix}{\marginpar{FIX}}
\newcommand{\new}{\marginpar{NEW}}

% ICLR 正文硬上限 9 页，浮动体间距默认偏松，这里收紧以腾出正文空间。
\setlength{\textfloatsep}{7pt plus 2pt minus 2pt}
\setlength{\floatsep}{7pt plus 2pt minus 2pt}
\setlength{\intextsep}{7pt plus 2pt minus 2pt}
% 正文有 6 张表 + 5 张图，默认每页浮动体上限会让它们堆积到最后一页。
% 放宽每页浮动体配额，让它们均匀分布在第 1-9 页内。
\setcounter{topnumber}{4}
\setcounter{bottomnumber}{3}
\setcounter{totalnumber}{5}
\renewcommand{\topfraction}{0.92}
\renewcommand{\bottomfraction}{0.65}
\renewcommand{\textfraction}{0.06}
\renewcommand{\floatpagefraction}{0.72}

\title{\introact{}: Selective Input Governance for \\
Frozen Time-Series Foundation Models}

\author{Anonymous Authors\\Paper under double-blind review}

%\iclrfinalcopy
\begin{document}

\maketitle

\begin{abstract}
Time-series foundation models are deployed without task-specific parameter updates, and the
context they receive at forecast time is often incomplete. Existing incomplete-input
forecasting methods reconstruct the unavailable information or train architectures that
tolerate missingness. We study a different question for a frozen forecaster. Given several
admissible ways of handling an incomplete context, which one should be applied to the current
request, and should any be applied at all. We introduce \introact{}, which replays historical
windows as pseudo-deployments to measure the realised forecasting utility of each admissible
action, estimates that utility locally for the action under consideration, and executes an
action only when a conservative utility score is positive. The reference action is part of the
catalog, so the method keeps the input unchanged whenever the evidence does not support a
switch, and every returned forecast comes from one input version that was actually executed.
Across \ph{N_DATASETS} sources, four missingness patterns, and two horizons, with two frozen
backbones used for development and a third family held out from method design,
\introact{} reaches \ph{MAIN_OVERALL_MASE} source-macro MASE against \ph{STRONG_BASELINE} at an
intervention rate of \ph{MAIN_INTERVENTION_RATE}, and the interventions it executes keep a
conditional harmful rate of \ph{HARM_HIR_OURS}. \ph{ABSTRACT_CLOSING}.
\end{abstract}

% ===========================================================================
\section{Introduction}
\label{sec:intro}

Time-series foundation models (\tsfm{}s) now produce competitive zero-shot forecasts without
task-specific retraining~\citep{ansari2024chronos,chronosbolt2025,das2024timesfm,woo2024moirai,
chronos2_2025}, which makes them attractive as fixed components of deployed forecasting
systems. The context available at the forecast origin is often incomplete. A sensor stops
reporting for a stretch of time, a channel delivers a variable several steps late, or only
part of an acquisition batch has arrived~\citep{chen2024bitgraph,peng2025s4m}. We study
deployment protocols in which an observation that has arrived and is valid must not be
overwritten, so the system may act only on positions where the data is missing.

The established response to missing data is imputation. Reconstruction-oriented methods
recover the missing entries as accurately as possible and are scored on reconstruction error.
BRITS uses bidirectional recurrent dynamics~\citep{cao2018brits}, CSDI models the conditional
distribution of the missing values with a score-based diffusion
process~\citep{tashiro2021csdi}, and T1~\citep{park2026t1} transfers cross-variable
information to improve multivariate recovery. A second line handles missingness inside the
forecasting model, where BiTGraph couples biased temporal convolution with a graph over
variables~\citep{chen2024bitgraph} and S4M folds the missing pattern into a structured
state-space model~\citep{peng2025s4m}. Both train or redesign the predictor.

A third line moves the objective from recovery to downstream utility. Task-oriented imputation
evaluation estimates how imputed values at different time steps affect a downstream
forecasting model and combines imputation strategies according to the estimated
gain~\citep{wang2024taskoriented}. TOI-VSF trains the imputation module to support the
forecasting objective~\citep{hao2025toivsf}, GIMCC adds causal consistency between the
generated values and the target~\citep{hao2025gimcc}, and SRDI imputes diffusively in a space
resilient to the shift induced by variable absence~\citep{xu2026srdi}. TATO adapts a frozen
\tsfm{} to heterogeneous domains through transformation pipelines selected from historical
task performance~\citep{qiu2026tato}. Each result describes a strategy or a transformation
family estimated over a corpus.

We study the decision that remains once the model is deployed. The forecaster is frozen, the
future target is unavailable, and the system must choose, for one specific incomplete request,
which admissible input intervention to execute or whether to keep the input unchanged. Three
properties make this choice non-trivial. The realised utility of an intervention depends on
the future target and is not observable online, which motivates replaying historical windows
as pseudo-deployments and recording what each admissible action actually did. The best action
changes with the request and with the backbone, so evidence has to be retrieved locally rather
than averaged over a corpus. Local evidence is finite and can conflict, and an action that
looks favourable on average can still increase the loss of the request at hand, so execution
is made conditional on a conservative utility score.

\paragraph{Contributions.}
\textbf{(1)~Selective input governance.} We formalise incomplete-input handling for a frozen
\tsfm{} as a per-request decision over a fixed catalog of admissible interventions, in which
the untouched input is one of the options, and we measure how strongly the best action varies
across requests and backbones (\S\ref{sec:exp-exists}). \textbf{(2)~Historical full-action
replay with local utility estimation.} We build a replay bank by re-enacting the registered
missingness protocol on TRAIN history and executing every catalog action with the frozen
forecaster, so the stored supervision is a realised forecasting utility, and the decision rule
returns the reference action whenever no action clears a conservative score
(\S\ref{sec:method}, \S\ref{sec:exp-harm}). \textbf{(3)~Evidence on harm and cost.} We
compare against reconstruction-based, task-oriented, and data-side adaptation baselines on
eight sources, three frozen \tsfm{} families, four missingness patterns, and three severity
levels, and we report forecasting accuracy, harmful interventions, and deployment cost in the
same protocol (\S\ref{sec:exp-main}, \S\ref{sec:exp-robust}, \S\ref{sec:exp-ablation}).

\begin{figure}[t]
  \centering
  \includegraphics[width=0.70\textwidth]{figure/fig1_selective_governance.pdf}
  \caption{Selective input governance for a frozen forecaster. \textbf{(a)}~Measured action
  utility on the evaluation episodes, drawn from the recorded statistics rather than from
  hand-picked examples. The oracle-best action is not constant across requests, so a single
  fixed intervention leaves part of the available improvement unused. \textbf{(b)}~Historical
  full-action replay. TRAIN windows are replayed under the registered missingness protocol,
  every catalog action is executed with the frozen forecaster, and the realised utilities are
  stored in the replay bank. \textbf{(c)}~Online decision. The reference forecast is computed,
  each admissible action is scored from the same-action neighbourhood of its action-conditioned
  state, and the best action is executed only when its conservative score is positive.}
  \label{fig:concept}
\end{figure}

% ===========================================================================
\section{Related Work}
\label{sec:related}

\subsection{Incomplete Time-Series Modeling}
\label{sec:related-incomplete}
Reconstruction-based imputation is the older of two routes. BRITS~\citep{cao2018brits}
performs bidirectional recurrent imputation, CSDI~\citep{tashiro2021csdi} replaces the
deterministic mapping with a conditional score-based diffusion model,
GAIN~\citep{yoon2018gain} and SAITS~\citep{du2023saits} explore adversarial and self-attention
variants of the same objective, and T1~\citep{park2026t1} improves multivariate recovery under
heavy and structurally different missingness. Their primary comparisons are centred on
reconstruction accuracy.

The second route forecasts from partially observed inputs without a separate imputation stage.
BiTGraph~\citep{chen2024bitgraph} points out that imputing first and forecasting afterwards
accumulates error, and handles the missing pattern inside a biased temporal convolutional
network with a variable graph. S4M~\citep{peng2025s4m} integrates missing-data handling into a
structured state-space architecture, and ChannelTokenFormer~\citep{jang2026channeltokenformer}
addresses inter-channel dependency, asynchronous sampling, and test-time missing blocks
jointly. Methods on this route are complete forecasting models whose parameters are part of the
contribution.

The third route moves the objective from recovery to downstream utility. Task-oriented
imputation evaluation~\citep{wang2024taskoriented} estimates how missing or imputed labels at
individual time steps affect the downstream forecasting model and combines imputation methods
according to the estimated gain. TOI-VSF~\citep{hao2025toivsf} carries the objective into
\emph{variable subset forecasting}, where whole variables are absent at inference.
GIMCC~\citep{hao2025gimcc} adds multi-level causal consistency between the generated values and
the target, VIDA~\citep{liang2025vida} reframes variable subset forecasting as cross-domain
knowledge transfer, and SRDI~\citep{xu2026srdi} imputes diffusively in a space resilient to the
shift induced by variable absence.

All three routes either recover the missing values, train a predictor that tolerates them, or
estimate a training-time property of an imputation method. \introact{} keeps the predictor
fixed and addresses the deployment stage, where the request has already arrived, the backbone
is already frozen, the future target is unavailable, and the system selects one input version
from a finite catalog or keeps the reference input.

\subsection{Data-Side Adaptation for Frozen \tsfm{}s}
\label{sec:related-dataside}
TATO~\citep{qiu2026tato} shows that a frozen \tsfm{} can be adapted to heterogeneous domains
without parameter updates. It searches a transformation pipeline, including context slicing,
scale normalisation, and outlier correction, and selects the pipeline from historical task
performance. The unit of adaptation is a target domain. TS-ICL~\citep{tsicl2026} is a
foundation model that natively supports imputation, irregular observations, and partially
observed look-back windows. We use it here only as the in-context regression backend of two
catalog actions.

The distinction from \introact{} is the unit of decision. TATO and TS-ICL change what the
model sees or what the model is. \introact{} leaves both untouched and decides, for one
request at a time, whether to apply an admissible input intervention.

\subsection{Selective Decision Making}
\label{sec:related-selective}
Selective prediction abstains on part of the input in order to trade coverage for lower
risk~\citep{chow1970optimum,vovk2005algorithmic,geifman2017selective,lei2018distribution}.
Learning to defer instead routes part of the input to a human or to another decision
maker~\citep{madras2018predict,mozannar2020consistent}. Both mechanisms change who or what
produces the answer. \introact{} does neither. The same frozen forecaster returns a forecast
in every case, and what is selected is whether the input is modified first.

The decision layer is also not a policy-learning problem over logged data. Historical replay
executes every action in the finite catalog on the same window and evaluates each against a
realised future, so the bank holds full action outcomes rather than bandit feedback from a
logging policy, and no propensity correction is required. Dynamic feature
selection~\citep{he2024dime} weighs the value of additional information against its cost in a
related sequential-acquisition setting. We make no information-theoretic optimality claim and
report the cost accounting separately (Appendix~\ref{app:calls}).

% ===========================================================================
\section{Problem Formulation and Method}
\label{sec:method}

\begin{figure}[t]
  \centering
  \includegraphics[width=0.48\textwidth]{figure/fig2_introact_architecture.pdf}
  \caption{Architecture of \introact{}. \textbf{Offline (top)}: historical TRAIN windows are
  turned into incomplete contexts by injecting the registered missingness protocol, every
  catalog action is executed with the frozen forecaster, and the realised utilities populate
  the replay bank. This is the only place where a future is read. \textbf{Online (bottom)}: a
  request exposes an incomplete context only. The reference forecast is computed, an
  action-conditioned state $z_a$ is formed for each admissible action, and each state is scored
  against the frozen bank. One action is then executed with the same frozen forecaster, and
  neither a future target nor a non-selected candidate's forecast is accessed online.}
  \label{fig:arch}
\end{figure}

\subsection{Problem Setup}
\label{sec:method-problem}
Let $X \in \mathbb{R}^{C \times L}$ be the multivariate context at the forecast origin, with
observation mask $M \in \{0,1\}^{C \times L}$, and let $Y \in \mathbb{R}^{C \times H}$ be the
future target, which is never available to the decision procedure. Let $F$ denote the complete
frozen deployment inference pipeline, which takes an input version and returns a forecast
$p = F(X)$ and includes whatever preprocessing the backbone applies internally, so a call to
$F$ is exactly the call the deployed system makes.

Let $\calA = \{a_0, a_1, \dots, a_J\}$ be a fixed catalog of admissible interventions, with
$J$ interventions in addition to the reference. Each action $a$ materialises an admissible
input version $X^{a} = T_a(X, M)$. The reference action $a_0$ returns the input unchanged and
corresponds to no external intervention, and $\calA^{+} = \calA \setminus \{a_0\}$ collects
the actions that modify the input. Let $\loss$ be the forecasting loss, defined on the same
normalised scale as the primary evaluation metric so that replay utilities and reported errors
are directly comparable. Utility is measured against the reference, so for a realised target
$Y$,
\begin{equation}
  g_a \;=\; \loss\!\left(F(X^{a_0}),\, Y\right) \;-\; \loss\!\left(F(X^{a}),\, Y\right),
  \label{eq:utility}
\end{equation}
and the task is to choose $a^\star(X, M) \in \calA$ without access to $Y$, while keeping the
number of frozen-backbone calls small.

Three constraints define the deployment protocol. No observation that has arrived and is valid
may be overwritten, so every $X^{a}$ differs from $X^{a_0}$ only at positions where $M = 0$.
The future target is unavailable, so no realised loss is observable at deployment time. The
reference forecast is the only forecast that may be computed before the decision, which is a
call budget rather than a statement about the future: a request may spend one
forecasting-backbone call to obtain $p_0 = F(X^{a_0})$, and the forecasts of the remaining
candidates must not be inspected when the action is chosen.

\subsection{Historical Full-Action Replay}
\label{sec:method-replay}
The quantity that decides the action, $g_a$, depends on $Y$ and is not observable online, but
it is computable for a historical window. Take a TRAIN window $i$ whose forecast target $y_i$
is already part of the available history, inject the registered missingness protocol to obtain
$(X_i, M_i)$, materialise every action in the catalog, and execute the frozen pipeline on each
version. The realised utility of action $a$ on window $i$ is
\begin{equation}
  g_{i,a} \;=\; \loss\!\left(F(T_{a_0}(X_i, M_i)),\, y_i\right)
  \;-\; \loss\!\left(F(T_{a}(X_i, M_i)),\, y_i\right).
  \label{eq:replay-gain}
\end{equation}
This is a measured outcome of an action the deployed system would not otherwise have taken,
recovered by re-executing a historical window rather than by modelling it, and both terms are
evaluated against the same realised future.

The state recorded with each outcome is action conditioned. For window $i$ and action $a$,
\begin{equation}
  z_{i,a} \;=\; \phi\!\left(X_i,\, M_i,\, X_i^{a},\, p_i^{0}\right),
  \qquad
  p_i^{0} = F(X_i^{a_0}),
  \label{eq:state}
\end{equation}
where $X_i^{a} = T_a(X_i, M_i)$. The state may use the reference forecast $p_i^0$, which the
protocol already permits, and it never uses the forecast of a non-selected candidate. The
replay bank is
\begin{equation}
  \calB \;=\; \bigl\{\, (z_{i,a},\, a,\, g_{i,a}) \,\bigr\},
  \label{eq:replay-bank}
\end{equation}
built once from TRAIN history and then frozen. Deployment reads it and never writes to it.

The realised utility in \eqref{eq:replay-gain} is measured over the forecast horizon, so
$g_{i,a}$ at $H=96$ and at $H=192$ are different quantities and must not share a
neighbourhood. The bank is therefore indexed by the frozen backbone revision, by the source,
and by the horizon, and a request is scored only against the bank slice that matches its own
horizon. It also covers only the missingness protocol registered in advance. These properties
are load-bearing, and Appendix~\ref{app:replay} states what they imply for the scope of the
method.

The step from the bank to a decision rests on one empirical assumption, stated explicitly.
Requests whose action-conditioned states are close in the state space have similar action
utilities, so the utility of an action on a new request can be estimated from the same-action
neighbourhood of its state in the bank. This local transfer assumption is not asserted outside
the registered grid, and the paper makes no claim about behaviour beyond the patterns and
severities that were replayed.

\subsection{Action-Conditioned Local Utility Estimation}
\label{sec:method-scoring}
The estimator is non-parametric and uses no learned encoder and no learned distance. The state
$z_a$ of \eqref{eq:state} carries four groups of features. Mask state describes how the request
is incomplete. Visible-context state is computed from the values that are present and
summarises level, trend, autocorrelation, seasonality, and recent volatility. Intervention
state describes the difference between $X^{a}$ and $X^{a_0}$. Reference-forecast state is
computed from $p_0 = F(X^{a_0})$ and summarises forecast level, range, trend, the jump from the
last visible observation, and roughness. The complete feature list, window lengths, and
normalisation statistics are given once in Appendix~\ref{app:features}, and continuous features
are standardised with statistics fitted on replay-fit parents only.

Retrieval is performed separately for each action. For action $a$, the neighbourhood
$\calN_a$ contains the $k$ bank records whose recorded action is $a$ and whose states are
closest to $z_a$ in standardised Euclidean distance. Each neighbour receives the weight
\begin{equation}
  w_i \;=\; \exp\!\left(-d_i / \tau\right),
  \qquad
  \tau \;=\; \operatorname{median}_{i \in \calN_a} d_i \;+\; \epsilon ,
  \label{eq:weights}
\end{equation}
so the bandwidth is set by the query neighbourhood itself and is not tuned. The local mean
utility, the local dispersion, and the effective sample size are
\begin{equation}
  \hat\mu_a \;=\; \frac{\sum_{i \in \calN_a} w_i\, g_{i,a}}{\sum_{i \in \calN_a} w_i},
  \qquad
  \hat\sigma_a^{2} \;=\; \frac{\sum_{i \in \calN_a} w_i\, (g_{i,a} - \hat\mu_a)^{2}}{\sum_{i \in \calN_a} w_i},
  \qquad
  n_{\mathrm{eff},a} \;=\; \frac{\left(\sum_{i \in \calN_a} w_i\right)^{2}}{\sum_{i \in \calN_a} w_i^{2}} .
  \label{eq:local-moments}
\end{equation}
The neighbourhood size $k$ is the only quantity searched here, over $\{8, 16, 32\}$ on the
TRAIN gate. The neighbourhood is defined by distance alone, and no threshold marks a request as
outside the replay support, so the estimator returns a value for every request it receives.
Section~\ref{sec:exp-robust} and Appendix~\ref{app:scope} treat the resulting scope as a
within-grid statement.

\subsection{Conservative Act-or-Keep Decision}
\label{sec:method-decision}
A positive local mean is not sufficient evidence for executing an action, because the
neighbourhood may be small or internally inconsistent. The decision therefore uses the
conservative score
\begin{equation}
  S_a \;=\; \hat\mu_a \;-\; \beta\, \frac{\hat\sigma_a}{\sqrt{n_{\mathrm{eff},a}}},
  \qquad a \in \calA^{+},
  \label{eq:score}
\end{equation}
where $\beta$ is a penalty strength, and
\begin{equation}
  \max_{a \in \calA^{+}} S_a > 0
  \;\Rightarrow\;
  a^{\star} = \arg\max_{a \in \calA^{+}} S_a ,
  \qquad\text{otherwise}\qquad
  a^{\star} = a_0 .
  \label{eq:decision}
\end{equation}
The score is a heuristic and we present it as one. It is not a calibrated lower bound on
$g_a$, it carries no coverage or risk guarantee, and $\beta = 1.64$ is one of the candidate
penalty strengths $\{0, 1, 1.64\}$ in the gate grid rather than a quantile of a sampling
distribution. The threshold of zero is not a tuned quantity either. Utility is defined relative
to the reference action in \eqref{eq:utility}, so a score at zero means that the evidence does
not distinguish the action from leaving the input unchanged.

The two hyperparameters $k$ and $\beta$ are selected on the TRAIN gate before any evaluation
run, under a criterion fixed in advance. Among the grid points whose harmful-intervention rate
on the gate stays at or below a cap of \ph{GATE_HARM_CAP}, we take the point with the best gate
MASE, so harm enters as a constraint rather than as a quantity to trade against accuracy. The
cap is a protocol constant frozen before any evaluation run and recorded with the resolved
configuration, so neither the cap nor the selected grid point is chosen on evaluation data. If
no grid point satisfies the cap, the configuration falls back to the most conservative grid
point, which is the one with the largest penalty strength. No evaluation data is used at any
step.

Two consequences of \eqref{eq:decision} matter. The reference action is always available, so
the rule always returns a forecast. This is a statement about availability only and does not
mean that the reference forecast is accurate or that returning it is safe. The rule also has
no way to detect that a request lies outside the replay support, so it applies to the
registered missingness grid and the results should not be read as out-of-grid generalisation.

\subsection{Deployment Procedure and Cost}
\label{sec:method-deploy}
A request is served in four steps. Read the visible context and mask and execute the reference
action to obtain $p_0$ and the reference-forecast state, then materialise the candidate input
versions and form the action-conditioned states $z_a$. Retrieve the same-action neighbourhoods
from the frozen bank and evaluate $S_a$ for every $a \in \calA^{+}$, then execute $a^{\star}$
and return the resulting forecast, which is the reference forecast when $a^{\star} = a_0$.

Four cost components are kept separate and reported individually in Appendix~\ref{app:calls},
namely the one-off offline bank construction, candidate materialisation, the
forecasting-backbone call, and retrieval overhead. A request costs one backbone call for the
reference action plus one more if an intervention is executed. The method does not ensemble
forecast outputs, does not correct them, does not fit a residual model, and does not update any
parameter of $F$.

% ===========================================================================
\section{Experiments}
\label{sec:exp}

\subsection{Experimental Setup}
\label{sec:exp-setup}
The main matrix covers \ph{N_DATASETS} public forecasting sources (ETTh1, ETTh2, ETTm1, ETTm2,
Electricity, Exchange, Traffic, and Weather) and three frozen \tsfm{} families:
Chronos-Bolt~\citep{chronosbolt2025}, TimesFM-2.5~\citep{das2024timesfm}, and
Chronos-2~\citep{chronos2_2025}. The first two are used for method development. Chronos-2
enters only after the configuration is frozen, with the identical algorithm and hyperparameters
and a replay bank rebuilt with Chronos-2 itself, so the held-out result is transfer of a frozen
configuration across backbone families rather than transfer of a replay bank across backbones.
Context length is $L = 512$ and the main horizons are $H \in \{96, 192\}$. Four deterministic
missingness patterns are frozen before any method runs, and every mask is derived by a
deterministic hash of source, parent, origin, horizon, pattern, and protocol seed. The main
comparison uses $10\%$ severity, and \S\ref{sec:exp-robust} adds $30\%$ and $50\%$ inside the
same registered grid. Appendix~\ref{app:data-protocol} defines the patterns, the derivation
rule, and the metric denominators.

The catalog is fixed at five admissible actions with \textsc{Keep} as the reference $a_0$, and
a complete and valid input is always mapped to \textsc{Keep}. Appendix~\ref{app:catalog}
defines each action and the status rules for unsupported, aliased, or failed executions.

The core comparison is fixed in advance and never reordered after results are seen. Four
published methods compose with each frozen backbone and therefore share one reference contract:
T1~\citep{park2026t1} covers reconstruction-based imputation,
TOI~\citep{wang2024taskoriented} covers task-oriented utility,
TATO~\citep{qiu2026tato} covers data-side adaptation of a frozen \tsfm{}, and
SRDI~\citep{xu2026srdi} covers missing-input repair.
ChannelTokenFormer~\citep{jang2026channeltokenformer} is a complete forecasting architecture
trained for the task rather than a data-side adapter, so it is reported as an external
forecasting reference and carries no per-backbone column. Three controls accompany the core
set: \textsc{Native Keep}, \textsc{Best Fixed}, and \textbf{R2-CART}. A \emph{Catalog Oracle}
that reads the future target is reported as a diagnostic reference and excluded from all ranks.
The severity matrix of \S\ref{sec:exp-robust} keeps the same frozen configuration and reuses a
reduced but type-complete set of seven methods, which omits SRDI so that the high-severity runs
stay bounded while every control type is still represented. Appendix~\ref{app:datasets}
documents the roster, the extended comparison, and the per-source tables.

Each source is split chronologically into TRAIN and TEST with an $L + H$ purge, and TRAIN is
partitioned by origin time into replay-fit, gate, and TRAIN-eval blocks. All design,
hyperparameter selection, feature normalisation, and baseline fitting use TRAIN only, and every
table is computed once on TEST after the configuration is frozen. MASE~\citep{hyndman2006mase}
is the primary metric and RMSSE the secondary one, aggregated source-macro with the parent as
the statistical unit and paired cluster bootstrap confidence intervals against each baseline.
Appendix~\ref{app:splits} records the split boundaries and the TEST manifest, and
Appendix~\ref{app:metric-conventions} gives the full metric and statistics conventions, and
Appendix~\ref{app:diagnostic-conventions} fixes the counting rules that the diagnostics rely
on, including tie-breaking for the oracle-best action and the treatment of zero-utility
interventions.

\subsection{Does Selective Intervention Need to Exist?}
\label{sec:exp-exists}
This section tests the two premises the method relies on. Both statistics are diagnostics
computed on TEST after the configuration is frozen, and neither is used for any selection.

The first premise is that the best action varies across requests. For each evaluation episode
we compute the oracle-best action as the minimiser of the realised loss over the catalog, and
we report the share of episodes on which each action is best together with the gap between the
best fixed catalog action and the catalog oracle. Figure~\ref{fig:concept}(a) shows the
distribution. \ph{MAIN_HET_SHARES}. \ph{MAIN_HET_GAP}. \ph{MAIN_FIXED_ORACLE_GAP}.

The second premise is that reconstruction quality and forecasting utility are not the same
criterion. For the methods that produce an explicit reconstruction, we compare the
reconstruction rank against the forecasting-utility rank within each parent and report winner
agreement and pairwise discordance, with a pooled Spearman correlation as a descriptive
summary. The comparison covers the core methods that emit an explicit estimate of the hidden
positions, and \introact{} is not placed in it, because it returns a forecast from an executed
input rather than a reconstruction of the missing values.
Figure~\ref{fig:why-selective}(a) shows the comparison, and Appendix~\ref{app:recutils} adds
the Reconstruction Oracle diagnostic. \ph{RANK_DISAGREE_STATS}.

\begin{figure}[t]
  \centering
  \includegraphics[width=0.70\textwidth]{figure/fig3_why_selective.pdf}
  \caption{Why the decision needs forecasting utility and why it has to be selective.
  \textbf{(a)}~Within-parent reconstruction rank against forecasting-utility rank. A lower
  reconstruction error does not imply a higher realised utility, so a reconstruction criterion
  is not a substitute for the utility the deployment decision is about. \textbf{(b)}~Risk
  against intervention rate. A low harmful rate is only meaningful together with the
  intervention rate at which it was obtained, because a method that almost never intervenes
  reaches a low harmful rate without being useful.}
  \label{fig:why-selective}
\end{figure}

The stratum-level breakdown of the same episodes is deferred to
Appendix~\ref{app:opportunity}, where the opportunity thresholds are also reported together
with their sensitivity to the boundary choice.

\subsection{Main Comparison}
\label{sec:exp-main}
\emph{How does \introact{} compare with published incomplete-input forecasting methods across
sources and frozen \tsfm{} backbones?}

Table~\ref{tab:main} reports the main comparison at $10\%$ severity over all eight sources,
three backbones, two horizons, and four patterns. \ph{MAIN_READING}.
\ph{MAIN_DELTA_CI}. \ph{MAIN_BACKBONE}. \ph{MAIN_BASELINE_SPREAD}. \ph{MAIN_R2CART}.
\ph{MAIN_RANK}. Per-source MASE, MAE, and RMSE are in Appendix~\ref{app:datasets}, and
deployment cost is reported separately in Appendix~\ref{app:calls}.

\begin{table}[t]
  \caption{Main comparison at $10\%$ missingness. MASE and RMSSE are source-macro averaged
  and lower is better everywhere. The first block holds the methods that compose with a frozen
  backbone and therefore share one reference contract, and these are the rows that carry a
  rank. \best{Bold} marks the best deployable method and \second{underline} the second best.
  ChannelTokenFormer is a complete forecasting architecture trained for the task rather than a
  data-side adapter, so it has no per-backbone column and is reported as an external
  forecasting reference outside the ranks. The oracle row reads the future target for diagnosis
  only and is excluded from the ranks. Chronos-2 is a held-out backbone whose configuration is
  not retuned.}
  \label{tab:main}
  \centering
  \small
  \setlength{\tabcolsep}{3.5pt}
\resizebox{\textwidth}{!}{\begin{tabular}{llccccc c}
    \toprule
    Method & Type & Bolt MASE $\downarrow$ & TimesFM MASE $\downarrow$
           & Chronos-2 MASE $\downarrow$ & Overall MASE $\downarrow$
           & Overall RMSSE $\downarrow$ & Avg.\ rank $\downarrow$ \\
    \midrule
    Native KEEP        & Frozen \tsfm{}          & \ph{KEEP_BOLT}   & \ph{KEEP_TF}   & \ph{KEEP_CH2}   & \ph{KEEP_OVR}   & \ph{KEEP_RMSSE}   & \ph{KEEP_RANK} \\
    Best Fixed         & Frozen \tsfm{}          & \ph{BF_BOLT}     & \ph{BF_TF}     & \ph{BF_CH2}     & \ph{BF_OVR}     & \ph{BF_RMSSE}     & \ph{BF_RANK} \\
    R2-CART            & Selector                & \ph{R2_BOLT}     & \ph{R2_TF}     & \ph{R2_CH2}     & \ph{R2_OVR}     & \ph{R2_RMSSE}     & \ph{R2_RANK} \\
    T1                 & Reconstruction          & \ph{T1_BOLT}     & \ph{T1_TF}     & \ph{T1_CH2}     & \ph{T1_OVR}     & \ph{T1_RMSSE}     & \ph{T1_RANK} \\
    TOI                & Task-oriented           & \ph{TOI_BOLT}    & \ph{TOI_TF}    & \ph{TOI_CH2}    & \ph{TOI_OVR}    & \ph{TOI_RMSSE}    & \ph{TOI_RANK} \\
    TATO               & Data-side adaptation    & \ph{TATO_BOLT}   & \ph{TATO_TF}   & \ph{TATO_CH2}   & \ph{TATO_OVR}   & \ph{TATO_RMSSE}   & \ph{TATO_RANK} \\
    SRDI               & Missing-input repair    & \ph{SRDI_BOLT}   & \ph{SRDI_TF}   & \ph{SRDI_CH2}   & \ph{SRDI_OVR}   & \ph{SRDI_RMSSE}   & \ph{SRDI_RANK} \\
    \midrule
    \rowcolor{bestgray}\introact{} & Frozen governance & \ph{OURS_BOLT} & \ph{OURS_TF} & \ph{OURS_CH2} & \ph{OURS_OVR} & \ph{OURS_RMSSE} & \ph{OURS_RANK} \\
    \midrule
    ChannelTokenFormer & External reference      & N/A\textsuperscript{*} & N/A\textsuperscript{*} & N/A\textsuperscript{*} & \ph{CTF_OVR} & \ph{CTF_RMSSE} & -- \\
    Catalog Oracle     & Diagnostic              & \oracle{\ph{ORACLE_BOLT}} & \oracle{\ph{ORACLE_TF}} & \oracle{\ph{ORACLE_CH2}} & \oracle{\ph{ORACLE_OVR}} & \oracle{\ph{ORACLE_RMSSE}} & \oracle{--} \\
    \bottomrule
  \end{tabular}}
  \vspace{-2pt}
  {\footnotesize\textsuperscript{*}Not applicable: ChannelTokenFormer is a trained forecasting
  architecture rather than a data-side adapter over a frozen backbone. It has no reference
  action on a frozen backbone, so it is excluded from the ranks rather than ranked on a
  different scale.}
\end{table}

\subsection{Does Selective Replay Reduce Harm?}
\label{sec:exp-harm}
\emph{Is the gain produced by choosing better interventions, or by intervening less often?}

Table~\ref{tab:harm} reports the selective-governance diagnostics at the default operating
point. An intervention is harmful when it increases the realised loss relative to leaving the
input unchanged, so the conditional harmful rate is conditioned on the requests where an
intervention was executed and must be read together with the intervention rate, because a
method that intervenes on almost nothing can report a low conditional rate without
contributing anything. Beneficial precision is the share of executed interventions that reduced
the loss, and missed opportunity is the share of episodes with a positive oracle opportunity on
which nothing was executed. \ph{HARM_READING}. \ph{HARM_IR}. \ph{HARM_HIR}. \ph{HARM_HL}.
\ph{HARM_BP}. \ph{HARM_MO}.

Figure~\ref{fig:why-selective}(b) traces the trade-off over the score grid by post-processing
the stored scores on TEST, with no threshold selected on TEST.
Appendix~\ref{app:governance-fig} reports the curve cell by cell. \ph{HARM_RISKCURVE}.

\begin{table}[t]
  \caption{Selective-governance diagnostics at the default operating point, $10\%$ severity.
  Intervention rate is the share of requests on which a non-reference action was executed.
  Conditional harmful rate is the share of executed interventions that increased the realised
  loss relative to \textsc{Keep}, and harmful loss is the mean increase over those requests.
  Beneficial precision is the share of executed interventions that reduced the loss, and missed
  opportunity is the share of positive-opportunity episodes on which nothing was executed.}
  \label{tab:harm}
  \centering
  \small
  \setlength{\tabcolsep}{3.5pt}
\resizebox{\textwidth}{!}{\begin{tabular}{lcccccc}
    \toprule
    Method & MASE $\downarrow$ & Intervention rate & Conditional HIR $\downarrow$ & Harmful loss $\downarrow$ & Beneficial precision $\uparrow$ & Missed opportunity $\downarrow$ \\
    \midrule
    Native KEEP        & \ph{HARM_MASE_KEEP}   & \ph{HARM_IR_KEEP}   & \ph{HARM_HIR_KEEP}   & \ph{HARM_HL_KEEP}   & \ph{HARM_BP_KEEP}   & \ph{HARM_MO_KEEP} \\
    Best Fixed         & \ph{HARM_MASE_BF}     & \ph{HARM_IR_BF}     & \ph{HARM_HIR_BF}     & \ph{HARM_HL_BF}     & \ph{HARM_BP_BF}     & \ph{HARM_MO_BF} \\
    R2-CART            & \ph{HARM_MASE_R2}     & \ph{HARM_IR_R2}     & \ph{HARM_HIR_R2}     & \ph{HARM_HL_R2}     & \ph{HARM_BP_R2}     & \ph{HARM_MO_R2} \\
    T1                 & \ph{HARM_MASE_T1}     & \ph{HARM_IR_T1}     & \ph{HARM_HIR_T1}     & \ph{HARM_HL_T1}     & \ph{HARM_BP_T1}     & \ph{HARM_MO_T1} \\
    TOI                & \ph{HARM_MASE_TOI}    & \ph{HARM_IR_TOI}    & \ph{HARM_HIR_TOI}    & \ph{HARM_HL_TOI}    & \ph{HARM_BP_TOI}    & \ph{HARM_MO_TOI} \\
    TATO               & \ph{HARM_MASE_TATO}   & \ph{HARM_IR_TATO}   & \ph{HARM_HIR_TATO}   & \ph{HARM_HL_TATO}   & \ph{HARM_BP_TATO}   & \ph{HARM_MO_TATO} \\
    SRDI               & \ph{HARM_MASE_SRDI}   & \ph{HARM_IR_SRDI}   & \ph{HARM_HIR_SRDI}   & \ph{HARM_HL_SRDI}   & \ph{HARM_BP_SRDI}   & \ph{HARM_MO_SRDI} \\
    \midrule
    \rowcolor{bestgray}\introact{} & \ph{HARM_MASE_OURS} & \ph{HARM_IR_OURS} & \ph{HARM_HIR_OURS} & \ph{HARM_HL_OURS} & \ph{HARM_BP_OURS} & \ph{HARM_MO_OURS} \\
    \bottomrule
  \end{tabular}}
  \vspace{-2pt}
  {\footnotesize Only methods that share the frozen-backbone reference contract appear here,
  because the harmful rate of an intervention is defined against the reference forecast of the
  same backbone. ChannelTokenFormer has no such reference and is therefore reported in
  Table~\ref{tab:main} only. Methods that always return a repaired input intervene on every
  incomplete request by construction, so their intervention rate is one and their harmful rate
  is unconditional.}
\end{table}

\subsection{Robustness and Transfer}
\label{sec:exp-robust}
\emph{Does one frozen configuration remain useful as missingness becomes more severe, and does
it carry over to a backbone family that was not used during development?}

Severities are $10\%$, $30\%$, and $50\%$, crossed with the four patterns and the three
backbones, and no method is retuned between levels. The levels lie inside the registered grid
that the replay bank samples from, so this is a within-grid robustness check and not an
unseen-severity stress test. The sweep uses a reduced but type-complete set of seven methods,
which omits SRDI so that the high-severity runs stay bounded, and the full core set is
compared at $10\%$ in Table~\ref{tab:main}. Table~\ref{tab:robust} reports the severity trend,
the worst pattern cell, and the held-out Chronos-2 backbone under an unchanged configuration.
\ph{ROBUST_READING}. \ph{ROBUST_30}. \ph{ROBUST_50}. \ph{ROBUST_WORST}. \ph{ROBUST_CH2}.
Appendix~\ref{app:severity-full} holds the complete per-severity tables and
Appendix~\ref{app:pattern} holds the per-pattern tables.

\begin{table}[t]
  \caption{Within-grid robustness to missingness severity, with no method retuned between
  levels. The seven rows cover every control type, and the same set is used at all three
  severities and in Appendix~\ref{app:severity-full}. The worst-pattern column is the largest
  degradation relative to \textsc{Native Keep} over all pattern and backbone cells at that
  severity. The Chronos-2 column is the held-out backbone under an unchanged configuration.}
  \label{tab:robust}
  \centering
  \small
  \setlength{\tabcolsep}{3pt}
\resizebox{\textwidth}{!}{\begin{tabular}{lccccc}
    \toprule
    Method & 10\% MASE $\downarrow$ & 30\% MASE $\downarrow$ & 50\% MASE $\downarrow$ & Worst pattern $\downarrow$ & Chronos-2 MASE $\downarrow$ \\
    \midrule
    Native KEEP        & \ph{RB_KEEP_10}    & \ph{RB_KEEP_30}    & \ph{RB_KEEP_50}    & \ph{RB_KEEP_WORST}    & \ph{RB_KEEP_CH2} \\
    Best Fixed         & \ph{RB_BF_10}      & \ph{RB_BF_30}      & \ph{RB_BF_50}      & \ph{RB_BF_WORST}      & \ph{RB_BF_CH2} \\
    R2-CART            & \ph{RB_R2_10}      & \ph{RB_R2_30}      & \ph{RB_R2_50}      & \ph{RB_R2_WORST}      & \ph{RB_R2_CH2} \\
    T1                 & \ph{RB_T1_10}      & \ph{RB_T1_30}      & \ph{RB_T1_50}      & \ph{RB_T1_WORST}      & \ph{RB_T1_CH2} \\
    TOI                & \ph{RB_TOI_10}     & \ph{RB_TOI_30}     & \ph{RB_TOI_50}     & \ph{RB_TOI_WORST}     & \ph{RB_TOI_CH2} \\
    TATO               & \ph{RB_TATO_10}    & \ph{RB_TATO_30}    & \ph{RB_TATO_50}    & \ph{RB_TATO_WORST}    & \ph{RB_TATO_CH2} \\
    \midrule
    \rowcolor{bestgray}\introact{} & \ph{RB_OURS_10} & \ph{RB_OURS_30} & \ph{RB_OURS_50} & \ph{RB_OURS_WORST} & \ph{RB_OURS_CH2} \\
    \bottomrule
  \end{tabular}}
\end{table}

\subsection{Ablation and Cost}
\label{sec:exp-ablation}
\emph{Which component carries the gain, and what does a request cost?}

Every ablation reuses the same catalog prediction cache, so the rows are recomputed selector
variants rather than new forecasting runs. A1 removes the dependence on the request, A2 the
intervention state, A3 the reference-forecast state, A4 the conservative penalty, and A5
replaces the non-parametric estimate with a parametric utility predictor, with the exact
definitions in Appendix~\ref{app:ablation-details}. Table~\ref{tab:ablation} reports accuracy,
harm, and cost together. \ph{ABL_READING}. \ph{ABL_A1}. \ph{ABL_A2}. \ph{ABL_A3}.
\ph{ABL_A4}. \ph{ABL_A5}. \ph{ABL_ORACLE}. \ph{COST_CALLS}. \ph{COST_LATENCY}.

\begin{table}[t]
  \caption{Core ablation and cost at $10\%$ severity. All rows read the same catalog
  prediction cache. Latency is end-to-end request latency and calls counts forecasting-backbone
  calls per request, with candidate materialisation reported in Appendix~\ref{app:calls}. The
  Reconstruction Oracle reads the hidden values and is excluded from every rank.}
  \label{tab:ablation}
  \centering
  \small
  \setlength{\tabcolsep}{3.5pt}
\resizebox{\textwidth}{!}{\begin{tabular}{lcccccc}
    \toprule
    Variant & MASE $\downarrow$ & Intervention rate & Conditional HIR $\downarrow$ & Harmful loss $\downarrow$ & Calls $\downarrow$ & Latency (ms) $\downarrow$ \\
    \midrule
    Full \introact{}          & \ph{ABL_FULL_MASE}   & \ph{ABL_FULL_IR}   & \ph{ABL_FULL_HIR}   & \ph{ABL_FULL_HL}   & \ph{ABL_FULL_CALLS}   & \ph{ABL_FULL_LAT} \\
    A1 Global utility          & \ph{ABL_A1_MASE}     & \ph{ABL_A1_IR}     & \ph{ABL_A1_HIR}     & \ph{ABL_A1_HL}     & \ph{ABL_A1_CALLS}     & \ph{ABL_A1_LAT} \\
    A2 w/o intervention state  & \ph{ABL_A2_MASE}     & \ph{ABL_A2_IR}     & \ph{ABL_A2_HIR}     & \ph{ABL_A2_HL}     & \ph{ABL_A2_CALLS}     & \ph{ABL_A2_LAT} \\
    A3 w/o reference-forecast state & \ph{ABL_A3_MASE} & \ph{ABL_A3_IR}     & \ph{ABL_A3_HIR}     & \ph{ABL_A3_HL}     & \ph{ABL_A3_CALLS}     & \ph{ABL_A3_LAT} \\
    A4 w/o conservative gate   & \ph{ABL_A4_MASE}     & \ph{ABL_A4_IR}     & \ph{ABL_A4_HIR}     & \ph{ABL_A4_HL}     & \ph{ABL_A4_CALLS}     & \ph{ABL_A4_LAT} \\
    A5 Parametric utility predictor & \ph{ABL_A5_MASE} & \ph{ABL_A5_IR}    & \ph{ABL_A5_HIR}     & \ph{ABL_A5_HL}     & \ph{ABL_A5_CALLS}     & \ph{ABL_A5_LAT} \\
    \midrule
    Reconstruction Oracle (diagnostic) & \oracle{\ph{ABL_ORACLE_MASE}} & \oracle{\ph{ABL_ORACLE_IR}} & \oracle{\ph{ABL_ORACLE_HIR}} & \oracle{\ph{ABL_ORACLE_HL}} & \oracle{--} & \oracle{--} \\
    \bottomrule
  \end{tabular}}
\end{table}

% ===========================================================================
\section{Conclusion}
\label{sec:conclusion}
We studied incomplete-input handling for frozen \tsfm{}s as a per-request decision over a
fixed catalog of admissible interventions. The central question is whether an intervention
improves forecasting for the current request relative to leaving the input unchanged. That
answer cannot be read off the future, which is unavailable, nor off reconstruction quality,
which is not the same criterion as realised forecasting utility. \introact{} answers it from
history by replaying TRAIN windows as pseudo-deployments, executing every catalog action with
the frozen forecaster, and storing the realised utility of each action with an
action-conditioned state. A request is then scored from the same-action neighbourhood of its
state, and an action is executed only when a conservative utility score is positive.

The evaluation supports three claims. \ph{CONCL_MAIN}. \ph{CONCL_HARM}. \ph{CONCL_TRANSFER}.
The evidence also has a clear boundary. All missingness here comes from controlled deletion
under a registered protocol, the replay bank covers only the patterns and severities that were
replayed, and the method cannot detect that a request lies outside that grid. The results
therefore support selective intervention under the evaluated controlled-missingness protocols
and frozen backbones, and generalisation to unseen mechanisms and to naturally incomplete
streams remains open. \ph{CONCLUSION_CLOSING}.

\subsubsection*{AI use statement}
We used generative AI tools to assist with literature review, prose editing, code drafting and
refactoring, figure preparation, and discussion of methodological and experimental choices.
All methodological decisions, experimental protocols, code changes, and reported results were
reviewed and approved by the authors. Generative AI was not used to fabricate or alter
experimental measurements, and no numerical result in this paper was generated by a language
model. Every reported number is produced by the recorded experiment pipeline described in
Appendix~\ref{app:reprodetails} and is traceable to a stored artefact. We take responsibility
for the final content of this work, including its text, claims, and artefacts.

\subsubsection*{Ethics statement}
This work studies input governance for frozen forecasting models on public benchmark datasets.
It does not involve human subjects, personal data, or deployment decisions with direct
societal consequences. We impose an integrity constraint under which an admissible
intervention cannot overwrite an observation that is marked as valid and available, because
silently replacing a real measurement would be a data integrity problem regardless of its
forecasting effect. We do not claim safety, robustness, or generality beyond the measured
conditions.

\subsubsection*{Reproducibility statement}
We provide the complete experimental protocol, the final TRAIN and TEST split with the
purging rule, the missingness generation procedure, the intervention catalog, the state
feature definitions, the hyperparameter search spaces, the failure policy, and the statistical
analysis in Appendix~\ref{app:data-protocol} and Appendix~\ref{app:splits}.
Appendix~\ref{app:reprodetails} specifies the result-record schema and the traceability checks.
All reported tables are generated from retained raw predictions and versioned configurations.

\bibliography{references}
\bibliographystyle{iclr2027_conference}

% ===========================================================================
% ===========================================================================
% 附录（由正文下沉 + 完整化）。所有未运行数字一律 \ph{KEY}，禁止编造。
% ===========================================================================



\appendix

\newpage



\section*{Appendix map}
\label{app:map}
Table~\ref{tab:app-map} lists what each appendix section contains. The appendix is organised
by evidence type rather than by individual analysis. Section~A fixes the protocol, the final
evaluation split, and the method details. Sections~B--D carry the evidence behind the
main-text tables. Section~E carries the ablations, the sensitivity grid, and the
mask-realisation stability check, and the replay-size study of our own bank sits in
Section~A with the rest of the bank construction. Sections~F and~G record cost and
reproducibility. A final unnumbered section holds exploratory material that does not support
any claim in the main text.

\begin{table}[h]
  \caption{Appendix map. Sections are lettered in order of appearance. The final
  supplementary section is unnumbered and holds no load-bearing evidence.}
  \label{tab:app-map}
  \centering
  \small
  \setlength{\tabcolsep}{4pt}
\resizebox{\textwidth}{!}{\begin{tabular}{lll}
    \toprule
    Section & Contents & Question it answers \\
    \midrule
    A. Protocol and method details & final TRAIN and TEST split with purging, scope and assumptions, the action catalog, the state feature definitions, the registered missingness patterns, the metric conventions, the diagnostic counting rules, the complete-input contract & what exactly was fixed before the evaluation ran, and which claims are out of scope? \\
    B. Full per-source and per-pattern results & per-source tables for every method, opportunity strata, per-pattern tables, extended comparison & does the aggregate hold source by source, stratum by stratum, and pattern by pattern? \\
    C. Full severity results & $10/30/50\%$ tables and the severity trend & does one frozen configuration hold across the registered severity grid? \\
    D. Reconstruction and utility diagnostics & within-parent rank comparison, Reconstruction Oracle, governance diagnostics & is reconstruction quality a usable substitute for forecasting utility? \\
    E. Ablations and stability & the full ablation vector, the $k$ and $\beta$ gate grid, and mask-realisation stability & which component carries the gain, and is the result an artefact of one mask? \\
    F. Cost, failures, and memory & latency, call counts, candidate materialisation, failures, timeouts, cold start & what does a request actually cost? \\
    G. Reproducibility and artefact schema & record schema, revisions, seeds, environment, manifests & can every reported number be traced back to a stored artefact? \\
    Supplementary (unnumbered) & the preceding development cycle & what was tried and rejected before the present formulation? \\
    \bottomrule
  \end{tabular}}
\end{table}



\section{Protocol and Method Details}
\label{app:protocol}
\subsection{Scope and Assumptions}
\label{app:scope}
This section states what the replay bank is, what it is not, and which claims the paper
therefore does and does not make.

The bank is a discrete grid over four patterns, three severities, and the registered sources,
and it is not a continuous model of missingness. If a request falls outside that grid, which
includes a different missingness mechanism, a different missing-rate regime, or an
unrepresented source family, the estimator still returns a score, but the retrieved
historical utilities need not transfer. The method has no guarantee outside the registered
replay support, because the neighbourhood is defined by distance alone and no threshold marks
a request as unsupported. We therefore read \S\ref{sec:exp-robust} as a within-grid check
rather than as out-of-grid generalisation, and results in this paper should not be interpreted
as out-of-grid generalisation.

Three further assumptions are worth stating explicitly. First, the action catalog is closed and
given. We do not search for new interventions, and a wider catalog would change every number
in this paper without changing the decision problem. Second, the missingness mechanism that
deployment encounters is assumed to be the one that the replay injects. A request produced by
a different mechanism is handled by the same estimator but lies outside the scope of the
reported results. Third, the frozen forecaster is assumed deterministic and revision pinned.
Every replay record is tied to the exact model revision that produced it, and a model update
invalidates the bank rather than silently degrading it.

Four claims are therefore not made anywhere in this paper. We do not claim out-of-grid or
out-of-distribution missingness generalisation. We do not claim that the method can detect
that historical evidence is insufficient. We do not claim any coverage or risk-control
guarantee for the conservative score. We do not claim that the results extend to naturally
incomplete deployment streams, since every missingness realisation in the main matrix comes
from controlled deletion under a registered protocol.


\subsection{Full Method Details}
\label{app:method-details}
This section expands \S\ref{sec:method} with the closed intervention catalog, the canonical
definition of the state features, and the relation of the decision layer to off-policy
evaluation and selective prediction. The state features are defined here once, and the main
text refers to this section rather than restating them.

\subsubsection{Intervention Catalog}
\label{app:catalog}
Table~\ref{tab:app-catalog} defines the five admissible actions. The catalog is closed, so no
action outside this set is ever executed, and no action may overwrite an observation that has
arrived and is valid.

\begin{table}[h]
  \caption{The five admissible interventions. Observed entries are the entries that are
  present and valid in the reference input $X^{a_0}$, and an action that modified them would
  violate the integrity constraint of \S\ref{sec:method-problem}.}
  \label{tab:app-catalog}
  \centering
  \small
  \setlength{\tabcolsep}{4pt}
\resizebox{\textwidth}{!}{\begin{tabular}{llcc}
    \toprule
    Action & Input version produced & May modify observed entries? & Extra backbone calls \\
    \midrule
    \textsc{Keep} $a_0$   & the incomplete context, unchanged & reference action & 0 \\
    \textsc{Ffill}        & forward-fill from the last valid value & no & 1 \\
    \textsc{Single TS-ICL}& univariate in-context regression on the target channel & no & 1 \\
    \textsc{Multi TS-ICL} & multivariate in-context regression using visible channels & no & 1 \\
    \textsc{Context Ridge}& ridge reconstruction from the visible context & no & 1 \\
    \bottomrule
  \end{tabular}}
\end{table}

A complete and valid input is always mapped to \textsc{Keep} and the reference forecast is
returned unchanged. Unsupported, aliased, and failed actions are recorded with an explicit
status, and they are never counted as successes. The required behaviour of the system in each
input regime is listed in \S\ref{app:contract}.

\subsubsection{State Features}
\label{app:features}
The state is action conditioned and is written $z_a$ for action $a$, with
\begin{equation}
  z_a \;=\; \phi\!\left(X,\, M,\, X^{a},\, p_0\right),
  \qquad p_0 = F(X^{a_0}) .
  \label{eq:app-state}
\end{equation}
It contains no future target information. It uses the reference forecast $p_0$, which the
call budget of \S\ref{sec:method-problem} already permits, and it never uses the forecast of a
non-selected candidate. Table~\ref{tab:app-features} is the canonical definition of the four
groups, including the window length, the channel aggregation, and the missing-value handling
of each feature. Continuous features are standardised with statistics fitted on replay-fit
parents only, and the normalisation is frozen with the rest of the configuration.

\begin{table}[h]
  \caption{Canonical definition of the action-conditioned state. Every feature is computable
  from the request and the reference forecast alone. Retrieval uses the standardised Euclidean
  distance between states, with no learned encoder and no learned metric.}
  \label{tab:app-features}
  \centering
  \small
  \setlength{\tabcolsep}{4pt}
\resizebox{\textwidth}{!}{\begin{tabular}{@{}l p{0.24\textwidth} p{0.30\textwidth} p{0.18\textwidth} p{0.13\textwidth}@{}}
    \toprule
    Group & Feature & Definition and window & Channel handling & Scaling \\
    \midrule
    Mask state (per request) & missing ratio & share of positions with $M = 0$ over the full context & over all channels, then averaged & none \\
     & missing runs & number of maximal runs of $M = 0$ & over all channels, then averaged & z-scored on replay-fit \\
     & longest run & length of the longest run of $M = 0$ & over all channels, then averaged & z-scored on replay-fit \\
     & distance to origin & steps from the nearest missing position to the forecast origin & minimum over channels & z-scored on replay-fit \\
     & pattern indicator & one-hot over P1--P4 & request level & none \\
    \midrule
    Visible-context state (per request) & level & median of the visible values over the last $W_v = 168$ steps & per channel, then averaged over visible channels & z-scored on replay-fit \\
     & local trend & ordinary least squares slope over the last $W_v$ visible steps & per channel, then averaged & z-scored on replay-fit \\
     & lag-one autocorrelation & correlation of the visible series with its one-step shift & per channel, then averaged & z-scored on replay-fit \\
     & seasonality strength & share of spectral energy in the dominant non-zero period & per channel, then averaged & z-scored on replay-fit \\
     & recent volatility & standard deviation of first differences over the last $W_v$ steps & per channel, then averaged & z-scored on replay-fit \\
    \midrule
    Intervention state (per action $a$) & action identity & one-hot over the $J$ non-reference actions & request level & none \\
     & changed fraction & share of positions where $X^{a}$ differs from $X^{a_0}$ & over all channels, then averaged & none \\
     & maximum change & largest absolute difference between $X^{a}$ and $X^{a_0}$, scaled by the visible scale & over all channels, then averaged & z-scored on replay-fit \\
     & change near origin & share of the change that falls within $W_n = 24$ steps of the origin & over all channels, then averaged & none \\
    \midrule
    Reference-forecast state (per request) & level and range & mean and range of $p_0$ over the horizon & per channel, then averaged & z-scored on replay-fit \\
     & forecast trend & ordinary least squares slope of $p_0$ over the horizon & per channel, then averaged & z-scored on replay-fit \\
     & first-step jump & first step of $p_0$ minus the last visible value, scaled by the visible scale & per channel, then averaged & z-scored on replay-fit \\
     & roughness & mean absolute second difference of $p_0$ & per channel, then averaged & z-scored on replay-fit \\
    \bottomrule
  \end{tabular}}
\end{table}

\subsubsection{Relation to Off-Policy Evaluation and Selective Prediction}
\label{app:ope-positioning}
The decision layer is related to conservative policy selection in contextual
bandits~\citep{abbasiyadkori2011improved,lattimore2020bandit} and to off-policy
evaluation~\citep{strehl2010learning,dudik2011doubly,swaminathan2015counterfactual,
jiang2016doubly,thomas2016data}, and the difference that matters is the structure of the
supervision. Off-policy evaluation estimates the value of a target policy from data collected
by a logging policy, so only the outcome of the logged action is observed and a propensity
correction is required. Historical replay gives full action outcomes for the finite catalog,
because every action is executed on the same window and evaluated against the same realised
future. There is therefore no partial-feedback problem to correct, and there is no online
exploration: no action is ever executed at deployment in order to learn.

Selective prediction abstains on part of the input in order to trade coverage for lower
risk~\citep{chow1970optimum,vovk2005algorithmic,geifman2017selective,lei2018distribution}, and
learning to defer routes part of the input to a human or to another decision
maker~\citep{madras2018predict,mozannar2020consistent}. Both change who or what produces the
answer. \introact{} changes neither. The same frozen forecaster returns a forecast in every
case, and what is selected is whether the input is modified first. Dynamic feature
selection~\citep{he2024dime} is a related sequential-acquisition setting in which the value of
additional information is weighed against its cost. We make no information-theoretic
optimality claim and report the cost accounting separately (Appendix~\ref{app:calls}).

\subsection{Replay Construction}
\label{app:replay}
Two properties of the bank carry the scope of the method. It is indexed by the frozen backbone
revision and by the source, since a utility measured with one revision of one model does not
describe another and the state distribution differs across sources. The method therefore
requires labelled in-source history for the backbone it serves, which places it in the
frozen-backbone and history-supervised regime rather than in a zero-shot regime. The bank also
covers only the missingness protocol registered in advance, so it is a discrete grid over
patterns and severities rather than a model of missingness.

This section records the split rule and how the replay bank is built. Nothing here is a
result. It is the construction that the result tables are read against.

\subsubsection{Splits, Purging, and the Final Evaluation Set}
\label{app:splits}
The split is defined from the outside in. Each source is first split chronologically into
TRAIN and TEST with an $L + H$ purge, so no TRAIN context window overlaps a TEST target window
and the two blocks share no origin. All method design, hyperparameter selection, feature
normalisation, baseline fitting, and acceptance checks use TRAIN only. The final tables are
computed once on TEST after the configuration is frozen, and no TEST parent participates in
any selection, normalisation, or bank construction.

TRAIN is then partitioned at the parent level. Within each source the registered admissible
TRAIN parents are sorted by origin time and divided into \emph{replay-fit}, \emph{gate}, and
\emph{TRAIN-eval} at approximately $60/20/20$. The replay bank is built exclusively from
replay-fit parents. The gate selects $k$, $\beta$, and any baseline hyperparameter.
TRAIN-eval is the final internal acceptance check before the configuration is frozen, and it
is not an evaluation set for any reported number. Purging is enforced at $L + H$ inside TRAIN
as well, so no parent's context window overlaps the target window of another split.
Table~\ref{tab:app-splits} records the realised sizes, including the TEST block.

\begin{table}[h]
  \caption{Realised split sizes. Counts are of admissible parents and not of correlated mask
  variants, so a parent with several patterns or severities still contributes one unit to
  every statistic in the paper. TEST is the evaluation population for every table in the main
  text, and it is disjoint from all three TRAIN blocks.}
  \label{tab:app-splits}
  \centering
  \small
  \setlength{\tabcolsep}{4pt}
\resizebox{\textwidth}{!}{\begin{tabular}{lcccc}
    \toprule
    Split & Parents & Origins & Sources covered & Used for \\
    \midrule
    Replay-fit   & \ph{FIT_PARENTS}   & \ph{FIT_ORIGINS}   & \ph{FIT_SOURCES}   & replay bank, baseline training, feature normalisation \\
    Gate         & \ph{GATE_PARENTS}  & \ph{GATE_ORIGINS}  & \ph{GATE_SOURCES}  & selection of $k$, $\beta$, and baseline hyperparameters \\
    TRAIN-eval   & \ph{EVAL_PARENTS}  & \ph{EVAL_ORIGINS}  & \ph{EVAL_SOURCES}  & frozen-configuration acceptance check \\
    \midrule
    TEST         & \ph{TEST_PARENTS}  & \ph{TEST_ORIGINS}  & \ph{TEST_SOURCES}  & every table in the main text, once, after freezing \\
    \bottomrule
  \end{tabular}}
\end{table}

The replay bank is constructed as follows. For every replay-fit parent we take the historical
window $(X_i, y_i)$ whose future is already available history, inject the registered
missingness protocol to obtain the incomplete version $X_i^{a_0}$, materialise every action
$a \in \calA$ into $X_i^{a}$, execute the frozen pipeline on each materialised version, and
store the realised utility
\begin{equation}
  g_{i,a} \;=\; \loss\big(F(T_{a_0}(X_i, M_i)),\, y_i\big) - \loss\big(F(T_{a}(X_i, M_i)),\, y_i\big),
  \label{eq:app-gain}
\end{equation}
so that $g_{i,a} > 0$ means the action reduced the forecasting loss and $g_{i,a} < 0$ means it
increased it. The bank is
$\calB = \{(z_{i,a}, a, g_{i,a})\}$, where $z_{i,a}$ is the action-conditioned state defined
in Appendix~\ref{app:features}. Because a utility measured with one model revision does not
describe another, because the state distribution differs across sources, and because $g_{i,a}$
is a loss over the forecast horizon and therefore a different quantity at $H=96$ and at
$H=192$, the bank is built per frozen backbone revision, per source, and per horizon. A request
is scored only against the slice that matches its own horizon, and no horizon indicator is
needed in the state because the horizon is fixed by the slice. Pseudo-deployment severities are drawn from
$\{10\%, 30\%, 50\%\}$ by deterministic hash of the same key that defines the mask, so the
bank covers the registered severity range without the method being retuned per severity.
Because the bank is built once from TRAIN and never rebuilt at deployment time, the offline
cost is paid once. Appendix~\ref{app:calls} separates it from the per-request online cost.

\subsubsection{Replay Bank Size}
\label{app:replaysize}
The replay bank is subsampled to $25\%$, $50\%$, and $100\%$ of its replay-fit support to
measure how much the method depends on the amount of historical evidence. This analysis is
not part of the main text and was not used to select any hyperparameter. The frozen
configuration is the one trained on the full bank. If performance at $25\%$ matches
performance at $100\%$, the honest reading is that the bank is already saturated at this
scale, not that the method is data-efficient in general.

\begin{table}[h]
  \caption{Replay-bank size ablation. The bank is subsampled by replay-fit parent, keeping
  every action and severity per retained parent, so the subsample changes the amount of
  history rather than its composition.}
  \label{tab:app-replaysize}
  \centering
  \small
  \setlength{\tabcolsep}{4pt}
  \begin{tabular}{lccccc}
    \toprule
    Bank fraction & MASE $\downarrow$ & Intervention Rate & HIR $\downarrow$ & Harmful Loss $\downarrow$ & Calls $\downarrow$ \\
    \midrule
    $25\%$  & \ph{RS25_MASE}  & \ph{RS25_IR}  & \ph{RS25_HIR}  & \ph{RS25_HL}  & \ph{RS25_CALLS} \\
    $50\%$  & \ph{RS50_MASE}  & \ph{RS50_IR}  & \ph{RS50_HIR}  & \ph{RS50_HL}  & \ph{RS50_CALLS} \\
    $100\%$ & \ph{RS100_MASE} & \ph{RS100_IR} & \ph{RS100_HIR} & \ph{RS100_HL} & \ph{RS100_CALLS} \\
    \bottomrule
  \end{tabular}
\end{table}

\subsection{Dataset and Missingness Protocol}
\label{app:data-protocol}
The missingness mechanism is frozen before any method is run, and every mask is derived
deterministically from a key that identifies the parent and the pattern. This subsection fixes
the four patterns, the derivation rule, and the metric conventions that depend on the data.

\subsubsection{Missingness Patterns}
\label{app:patterns}
Table~\ref{tab:app-patterns} defines the four deterministic patterns. Every mask is derived
by a deterministic hash of
\texttt{source + parent + origin + horizon + pattern + protocol\_seed} and is never moved
after forecasts are seen. The same key determines the pseudo-deployment severity that the
replay uses, which is drawn from $\{10\%, 30\%, 50\%\}$.

\begin{table}[h]
  \caption{The four registered missingness patterns. All four are applied to the same parents,
  so pattern effects are measured on a common parent set rather than on different subsets of
  the data.}
  \label{tab:app-patterns}
  \centering
  \small
  \setlength{\tabcolsep}{4pt}
\begin{tabular}{@{}l p{0.34\textwidth} p{0.40\textwidth}@{}}
    \toprule
    Pattern & What is removed & Why it is included \\
    \midrule
    P1 point missing    & scattered single positions & the classical random-missing regime \\
    P2 target block     & a contiguous block on the target variable & a sensor outage on the variable being forecast \\
    P3 shared block     & several visible variables over the same interval & a channel or acquisition outage affecting the whole panel \\
    P4 tail missing     & a block adjacent to the forecast origin & observations that have not arrived yet at request time \\
    \bottomrule
  \end{tabular}
\end{table}

\subsubsection{Metric Conventions}
\label{app:metric-conventions}
The statistical unit is the parent. Paired cluster bootstrap with \ph{N_BOOTSTRAP} resamples
over parent blocks gives 95\% confidence intervals on source-macro MASE differences against
each baseline, the source-macro aggregation is applied inside every resample, and Holm
correction~\citep{holm1979simple} is applied across the pre-declared baseline comparisons.
Average rank and won cells are descriptive summaries only, and the headline claim rests on
paired differences with confidence intervals. Mask variants of the same parent are never
treated as independent samples, and pattern, severity, and backbone breakdowns are secondary
analyses without cell-by-cell significance testing.

MASE~\citep{hyndman2006mase} is the primary metric and RMSSE is the secondary metric. Both
scale the error of a source by its own seasonal naive error, computed on the TRAIN portion of
that source and therefore never on the evaluation block. With $y_t$ the target channel and $m$
the registered seasonal period of the source,
\begin{equation}
  \mathrm{MASE} = \frac{\mathrm{MAE}}{\frac{1}{T-m}\sum_{t} \lvert y_t - y_{t-m} \rvert},
  \qquad
  \mathrm{RMSSE} = \sqrt{\frac{\mathrm{MSE}}{\frac{1}{T-m}\sum_{t} (y_t - y_{t-m})^2}},
  \label{eq:metrics}
\end{equation}
so the two metrics share the same seasonal naive difference series but not the same power of
it, and RMSSE is not the square root of MASE. The scaling constants are computed per channel
before any aggregation over channels, and the registered periods are listed in
Table~\ref{tab:app-seasonality}. A degenerate seasonal naive error, meaning a constant or
all-missing TRAIN series, leaves the metric undefined and is reported as missing rather than
as zero. Aggregation is source-macro, first over parents within a source and then equally over
the eight sources, and the same order is applied inside every bootstrap resample. MAE and RMSE
are reported per source and are never averaged across sources, because the sources have
different scales and a cross-source average of a raw error has no common unit.

\begin{table}[h]
  \caption{Per-source seasonal periods and context conventions used by the metric
  denominators. Context length is $L = 512$ for every source. The seasonal period is the
  sampling-frequency period registered before any run, not a tuned quantity.}
  \label{tab:app-seasonality}
  \centering
  \small
  \setlength{\tabcolsep}{4pt}
\resizebox{\textwidth}{!}{\begin{tabular}{lccc}
    \toprule
    Source & Seasonal period $m$ & Channels & Scaling \\
    \midrule
    ETTh1 / ETTh2 & 24 & 7 & TRAIN seasonal naive error of the target channel \\
    ETTm1 / ETTm2 & 96 & 7 & TRAIN seasonal naive error of the target channel \\
    Electricity & 24 & 321 & TRAIN seasonal naive error of the target channel \\
    Exchange & 5 & 8 & TRAIN seasonal naive error of the target channel \\
    Traffic & 24 & 862 & TRAIN seasonal naive error of the target channel \\
    Weather & 144 & 21 & TRAIN seasonal naive error of the target channel \\
    \bottomrule
  \end{tabular}}
\end{table}

\subsubsection{Diagnostic Conventions}
\label{app:diagnostic-conventions}
Four quantities in the diagnostics are counting rules rather than metrics, and each one is
fixed before any evaluation record is read, because a different convention would change the
reported numbers without changing a single forecast.

\emph{Oracle-best action.} The oracle-best action of an episode is the admissible action with
the smallest realised loss. Ties are broken by the frozen catalog order, so the first action in
that order wins and the reported shares of episodes on which each action is best sum to one. An
action that is unsupported, aliased, or failed is not admissible and is excluded from the
argmin rather than being counted as a loss.

\emph{Utility sign.} Utility is defined relative to the reference action in
\eqref{eq:utility}, so $g_{i,a} = 0$ means the action and the reference are indistinguishable on
that episode. An intervention counts as harmful only when $g_{i,a} < 0$ strictly and as
beneficial only when $g_{i,a} > 0$ strictly, so the zero cases enter neither the conditional
harmful rate nor beneficial precision and are reported as their own count.

\emph{Beneficial precision.} It is the number of executed interventions with strictly positive
utility divided by the number of executed interventions, and it is undefined when nothing was
executed. It is reported as missing rather than as zero in that case.

\emph{Opportunity strata.} The no-op boundary and the low-to-high boundary are the two
pre-declared quantiles of $\Delta_i$ computed on the gate split, computed once and applied
unchanged to TEST. The same two values partition every method, and no boundary is recomputed
per method. Both values are recorded with the resolved configuration.

\subsection{Complete-Input Contract}
\label{app:contract}
Under a complete and valid input, \introact{} must return \textsc{Keep}. No forecasting
improvement is claimed in this regime. The purpose of this subsection is to document the
governance contract rather than to report a result. Table~\ref{tab:app-complete} lists the
regimes and the required behaviour in each.

\begin{table}[h]
  \caption{Governance contract by input regime. The complete-input rows are contract
  statements and not accuracy results. The expected behaviour there is a no-op, and a non-no-op
  would be a defect rather than an improvement.}
  \label{tab:app-complete}
  \centering
  \small
  \setlength{\tabcolsep}{4pt}
\resizebox{\textwidth}{!}{\begin{tabular}{lll}
    \toprule
    Input regime & Required behaviour & Claim \\
    \midrule
    Complete and valid input            & return \textsc{Keep} and the reference forecast & no-op contract, no accuracy claim \\
    Complete input, malformed request   & reject with an explicit status & no claim \\
    Unsupported action requested        & reject the action, keep the catalog closed & no claim \\
    Incomplete input, no action scores above zero & return the reference action & the reference action is always available \\
    Incomplete input, one action scores above zero & execute exactly one action and return its forecast & evaluated in \S\ref{sec:exp-main} \\
    \bottomrule
  \end{tabular}}
\end{table}


\section{Full Per-Source and Per-Pattern Results}
\label{app:full-results}
\subsubsection{Per-Source Results}
\label{app:datasets}
Each table below reports one frozen backbone at $10\%$ severity, one horizon, and
all eight sources, with one row per method. MASE is reported in every cell so that a
reader can recompute the source-macro aggregate of Table~\ref{tab:main} directly from
these tables. The \emph{macro} column is the equal-weight average over the eight sources
and is the quantity the main table reports. Rows are generated from the recorded result
group rather than transcribed by hand, so the appendix and the main table cannot disagree.
MAE and RMSE are reported in the same tables in Appendix~\ref{app:cost-full} style and
are not averaged across sources.

\begin{table}[h]
  \caption{Per-source MASE on Bolt at $10\%$ severity, $H=96$. The macro column is the
  equal-weight average over the eight sources and is the quantity reported in
  Table~\ref{tab:main}. Lower is better everywhere.}
  \label{tab:app-src-bolt-96}
  \centering
  \small
  \setlength{\tabcolsep}{3pt}
\resizebox{\textwidth}{!}{\begin{tabular}{lcccccccc c}
    \toprule
    Method & ETTh1 & ETTh2 & ETTm1 & ETTm2 & Electricity & Exchange & Traffic & Weather & Macro $\downarrow$ \\
    \midrule
    Native KEEP & \ph{SRC_BOLT_KEEP_96_ETTH1} & \ph{SRC_BOLT_KEEP_96_ETTH2} & \ph{SRC_BOLT_KEEP_96_ETTM1} & \ph{SRC_BOLT_KEEP_96_ETTM2} & \ph{SRC_BOLT_KEEP_96_ELEC} & \ph{SRC_BOLT_KEEP_96_EXCH} & \ph{SRC_BOLT_KEEP_96_TRAF} & \ph{SRC_BOLT_KEEP_96_WEA} & \ph{SRC_BOLT_KEEP_96_MACRO} \\
    Best Fixed & \ph{SRC_BOLT_BF_96_ETTH1} & \ph{SRC_BOLT_BF_96_ETTH2} & \ph{SRC_BOLT_BF_96_ETTM1} & \ph{SRC_BOLT_BF_96_ETTM2} & \ph{SRC_BOLT_BF_96_ELEC} & \ph{SRC_BOLT_BF_96_EXCH} & \ph{SRC_BOLT_BF_96_TRAF} & \ph{SRC_BOLT_BF_96_WEA} & \ph{SRC_BOLT_BF_96_MACRO} \\
    R2-CART & \ph{SRC_BOLT_R2_96_ETTH1} & \ph{SRC_BOLT_R2_96_ETTH2} & \ph{SRC_BOLT_R2_96_ETTM1} & \ph{SRC_BOLT_R2_96_ETTM2} & \ph{SRC_BOLT_R2_96_ELEC} & \ph{SRC_BOLT_R2_96_EXCH} & \ph{SRC_BOLT_R2_96_TRAF} & \ph{SRC_BOLT_R2_96_WEA} & \ph{SRC_BOLT_R2_96_MACRO} \\
    T1 & \ph{SRC_BOLT_T1_96_ETTH1} & \ph{SRC_BOLT_T1_96_ETTH2} & \ph{SRC_BOLT_T1_96_ETTM1} & \ph{SRC_BOLT_T1_96_ETTM2} & \ph{SRC_BOLT_T1_96_ELEC} & \ph{SRC_BOLT_T1_96_EXCH} & \ph{SRC_BOLT_T1_96_TRAF} & \ph{SRC_BOLT_T1_96_WEA} & \ph{SRC_BOLT_T1_96_MACRO} \\
    TOI & \ph{SRC_BOLT_TOI_96_ETTH1} & \ph{SRC_BOLT_TOI_96_ETTH2} & \ph{SRC_BOLT_TOI_96_ETTM1} & \ph{SRC_BOLT_TOI_96_ETTM2} & \ph{SRC_BOLT_TOI_96_ELEC} & \ph{SRC_BOLT_TOI_96_EXCH} & \ph{SRC_BOLT_TOI_96_TRAF} & \ph{SRC_BOLT_TOI_96_WEA} & \ph{SRC_BOLT_TOI_96_MACRO} \\
    TATO & \ph{SRC_BOLT_TATO_96_ETTH1} & \ph{SRC_BOLT_TATO_96_ETTH2} & \ph{SRC_BOLT_TATO_96_ETTM1} & \ph{SRC_BOLT_TATO_96_ETTM2} & \ph{SRC_BOLT_TATO_96_ELEC} & \ph{SRC_BOLT_TATO_96_EXCH} & \ph{SRC_BOLT_TATO_96_TRAF} & \ph{SRC_BOLT_TATO_96_WEA} & \ph{SRC_BOLT_TATO_96_MACRO} \\
    SRDI & \ph{SRC_BOLT_SRDI_96_ETTH1} & \ph{SRC_BOLT_SRDI_96_ETTH2} & \ph{SRC_BOLT_SRDI_96_ETTM1} & \ph{SRC_BOLT_SRDI_96_ETTM2} & \ph{SRC_BOLT_SRDI_96_ELEC} & \ph{SRC_BOLT_SRDI_96_EXCH} & \ph{SRC_BOLT_SRDI_96_TRAF} & \ph{SRC_BOLT_SRDI_96_WEA} & \ph{SRC_BOLT_SRDI_96_MACRO} \\
    \midrule
    \rowcolor{bestgray}\introact{} & \ph{SRC_BOLT_OURS_96_ETTH1} & \ph{SRC_BOLT_OURS_96_ETTH2} & \ph{SRC_BOLT_OURS_96_ETTM1} & \ph{SRC_BOLT_OURS_96_ETTM2} & \ph{SRC_BOLT_OURS_96_ELEC} & \ph{SRC_BOLT_OURS_96_EXCH} & \ph{SRC_BOLT_OURS_96_TRAF} & \ph{SRC_BOLT_OURS_96_WEA} & \ph{SRC_BOLT_OURS_96_MACRO} \\
    \midrule
    Catalog Oracle & \oracle{\ph{SRC_BOLT_ORACLE_96_ETTH1}} & \oracle{\ph{SRC_BOLT_ORACLE_96_ETTH2}} & \oracle{\ph{SRC_BOLT_ORACLE_96_ETTM1}} & \oracle{\ph{SRC_BOLT_ORACLE_96_ETTM2}} & \oracle{\ph{SRC_BOLT_ORACLE_96_ELEC}} & \oracle{\ph{SRC_BOLT_ORACLE_96_EXCH}} & \oracle{\ph{SRC_BOLT_ORACLE_96_TRAF}} & \oracle{\ph{SRC_BOLT_ORACLE_96_WEA}} & \oracle{\ph{SRC_BOLT_ORACLE_96_MACRO}} \\
    \bottomrule
  \end{tabular}}
\end{table}

\begin{table}[h]
  \caption{Per-source MASE on Bolt at $10\%$ severity, $H=192$. The macro column is the
  equal-weight average over the eight sources and is the quantity reported in
  Table~\ref{tab:main}. Lower is better everywhere.}
  \label{tab:app-src-bolt-192}
  \centering
  \small
  \setlength{\tabcolsep}{3pt}
\resizebox{\textwidth}{!}{\begin{tabular}{lcccccccc c}
    \toprule
    Method & ETTh1 & ETTh2 & ETTm1 & ETTm2 & Electricity & Exchange & Traffic & Weather & Macro $\downarrow$ \\
    \midrule
    Native KEEP & \ph{SRC_BOLT_KEEP_192_ETTH1} & \ph{SRC_BOLT_KEEP_192_ETTH2} & \ph{SRC_BOLT_KEEP_192_ETTM1} & \ph{SRC_BOLT_KEEP_192_ETTM2} & \ph{SRC_BOLT_KEEP_192_ELEC} & \ph{SRC_BOLT_KEEP_192_EXCH} & \ph{SRC_BOLT_KEEP_192_TRAF} & \ph{SRC_BOLT_KEEP_192_WEA} & \ph{SRC_BOLT_KEEP_192_MACRO} \\
    Best Fixed & \ph{SRC_BOLT_BF_192_ETTH1} & \ph{SRC_BOLT_BF_192_ETTH2} & \ph{SRC_BOLT_BF_192_ETTM1} & \ph{SRC_BOLT_BF_192_ETTM2} & \ph{SRC_BOLT_BF_192_ELEC} & \ph{SRC_BOLT_BF_192_EXCH} & \ph{SRC_BOLT_BF_192_TRAF} & \ph{SRC_BOLT_BF_192_WEA} & \ph{SRC_BOLT_BF_192_MACRO} \\
    R2-CART & \ph{SRC_BOLT_R2_192_ETTH1} & \ph{SRC_BOLT_R2_192_ETTH2} & \ph{SRC_BOLT_R2_192_ETTM1} & \ph{SRC_BOLT_R2_192_ETTM2} & \ph{SRC_BOLT_R2_192_ELEC} & \ph{SRC_BOLT_R2_192_EXCH} & \ph{SRC_BOLT_R2_192_TRAF} & \ph{SRC_BOLT_R2_192_WEA} & \ph{SRC_BOLT_R2_192_MACRO} \\
    T1 & \ph{SRC_BOLT_T1_192_ETTH1} & \ph{SRC_BOLT_T1_192_ETTH2} & \ph{SRC_BOLT_T1_192_ETTM1} & \ph{SRC_BOLT_T1_192_ETTM2} & \ph{SRC_BOLT_T1_192_ELEC} & \ph{SRC_BOLT_T1_192_EXCH} & \ph{SRC_BOLT_T1_192_TRAF} & \ph{SRC_BOLT_T1_192_WEA} & \ph{SRC_BOLT_T1_192_MACRO} \\
    TOI & \ph{SRC_BOLT_TOI_192_ETTH1} & \ph{SRC_BOLT_TOI_192_ETTH2} & \ph{SRC_BOLT_TOI_192_ETTM1} & \ph{SRC_BOLT_TOI_192_ETTM2} & \ph{SRC_BOLT_TOI_192_ELEC} & \ph{SRC_BOLT_TOI_192_EXCH} & \ph{SRC_BOLT_TOI_192_TRAF} & \ph{SRC_BOLT_TOI_192_WEA} & \ph{SRC_BOLT_TOI_192_MACRO} \\
    TATO & \ph{SRC_BOLT_TATO_192_ETTH1} & \ph{SRC_BOLT_TATO_192_ETTH2} & \ph{SRC_BOLT_TATO_192_ETTM1} & \ph{SRC_BOLT_TATO_192_ETTM2} & \ph{SRC_BOLT_TATO_192_ELEC} & \ph{SRC_BOLT_TATO_192_EXCH} & \ph{SRC_BOLT_TATO_192_TRAF} & \ph{SRC_BOLT_TATO_192_WEA} & \ph{SRC_BOLT_TATO_192_MACRO} \\
    SRDI & \ph{SRC_BOLT_SRDI_192_ETTH1} & \ph{SRC_BOLT_SRDI_192_ETTH2} & \ph{SRC_BOLT_SRDI_192_ETTM1} & \ph{SRC_BOLT_SRDI_192_ETTM2} & \ph{SRC_BOLT_SRDI_192_ELEC} & \ph{SRC_BOLT_SRDI_192_EXCH} & \ph{SRC_BOLT_SRDI_192_TRAF} & \ph{SRC_BOLT_SRDI_192_WEA} & \ph{SRC_BOLT_SRDI_192_MACRO} \\
    \midrule
    \rowcolor{bestgray}\introact{} & \ph{SRC_BOLT_OURS_192_ETTH1} & \ph{SRC_BOLT_OURS_192_ETTH2} & \ph{SRC_BOLT_OURS_192_ETTM1} & \ph{SRC_BOLT_OURS_192_ETTM2} & \ph{SRC_BOLT_OURS_192_ELEC} & \ph{SRC_BOLT_OURS_192_EXCH} & \ph{SRC_BOLT_OURS_192_TRAF} & \ph{SRC_BOLT_OURS_192_WEA} & \ph{SRC_BOLT_OURS_192_MACRO} \\
    \midrule
    Catalog Oracle & \oracle{\ph{SRC_BOLT_ORACLE_192_ETTH1}} & \oracle{\ph{SRC_BOLT_ORACLE_192_ETTH2}} & \oracle{\ph{SRC_BOLT_ORACLE_192_ETTM1}} & \oracle{\ph{SRC_BOLT_ORACLE_192_ETTM2}} & \oracle{\ph{SRC_BOLT_ORACLE_192_ELEC}} & \oracle{\ph{SRC_BOLT_ORACLE_192_EXCH}} & \oracle{\ph{SRC_BOLT_ORACLE_192_TRAF}} & \oracle{\ph{SRC_BOLT_ORACLE_192_WEA}} & \oracle{\ph{SRC_BOLT_ORACLE_192_MACRO}} \\
    \bottomrule
  \end{tabular}}
\end{table}

\begin{table}[h]
  \caption{Per-source MASE on TimesFM at $10\%$ severity, $H=96$. The macro column is the
  equal-weight average over the eight sources and is the quantity reported in
  Table~\ref{tab:main}. Lower is better everywhere.}
  \label{tab:app-src-tf-96}
  \centering
  \small
  \setlength{\tabcolsep}{3pt}
\resizebox{\textwidth}{!}{\begin{tabular}{lcccccccc c}
    \toprule
    Method & ETTh1 & ETTh2 & ETTm1 & ETTm2 & Electricity & Exchange & Traffic & Weather & Macro $\downarrow$ \\
    \midrule
    Native KEEP & \ph{SRC_TF_KEEP_96_ETTH1} & \ph{SRC_TF_KEEP_96_ETTH2} & \ph{SRC_TF_KEEP_96_ETTM1} & \ph{SRC_TF_KEEP_96_ETTM2} & \ph{SRC_TF_KEEP_96_ELEC} & \ph{SRC_TF_KEEP_96_EXCH} & \ph{SRC_TF_KEEP_96_TRAF} & \ph{SRC_TF_KEEP_96_WEA} & \ph{SRC_TF_KEEP_96_MACRO} \\
    Best Fixed & \ph{SRC_TF_BF_96_ETTH1} & \ph{SRC_TF_BF_96_ETTH2} & \ph{SRC_TF_BF_96_ETTM1} & \ph{SRC_TF_BF_96_ETTM2} & \ph{SRC_TF_BF_96_ELEC} & \ph{SRC_TF_BF_96_EXCH} & \ph{SRC_TF_BF_96_TRAF} & \ph{SRC_TF_BF_96_WEA} & \ph{SRC_TF_BF_96_MACRO} \\
    R2-CART & \ph{SRC_TF_R2_96_ETTH1} & \ph{SRC_TF_R2_96_ETTH2} & \ph{SRC_TF_R2_96_ETTM1} & \ph{SRC_TF_R2_96_ETTM2} & \ph{SRC_TF_R2_96_ELEC} & \ph{SRC_TF_R2_96_EXCH} & \ph{SRC_TF_R2_96_TRAF} & \ph{SRC_TF_R2_96_WEA} & \ph{SRC_TF_R2_96_MACRO} \\
    T1 & \ph{SRC_TF_T1_96_ETTH1} & \ph{SRC_TF_T1_96_ETTH2} & \ph{SRC_TF_T1_96_ETTM1} & \ph{SRC_TF_T1_96_ETTM2} & \ph{SRC_TF_T1_96_ELEC} & \ph{SRC_TF_T1_96_EXCH} & \ph{SRC_TF_T1_96_TRAF} & \ph{SRC_TF_T1_96_WEA} & \ph{SRC_TF_T1_96_MACRO} \\
    TOI & \ph{SRC_TF_TOI_96_ETTH1} & \ph{SRC_TF_TOI_96_ETTH2} & \ph{SRC_TF_TOI_96_ETTM1} & \ph{SRC_TF_TOI_96_ETTM2} & \ph{SRC_TF_TOI_96_ELEC} & \ph{SRC_TF_TOI_96_EXCH} & \ph{SRC_TF_TOI_96_TRAF} & \ph{SRC_TF_TOI_96_WEA} & \ph{SRC_TF_TOI_96_MACRO} \\
    TATO & \ph{SRC_TF_TATO_96_ETTH1} & \ph{SRC_TF_TATO_96_ETTH2} & \ph{SRC_TF_TATO_96_ETTM1} & \ph{SRC_TF_TATO_96_ETTM2} & \ph{SRC_TF_TATO_96_ELEC} & \ph{SRC_TF_TATO_96_EXCH} & \ph{SRC_TF_TATO_96_TRAF} & \ph{SRC_TF_TATO_96_WEA} & \ph{SRC_TF_TATO_96_MACRO} \\
    SRDI & \ph{SRC_TF_SRDI_96_ETTH1} & \ph{SRC_TF_SRDI_96_ETTH2} & \ph{SRC_TF_SRDI_96_ETTM1} & \ph{SRC_TF_SRDI_96_ETTM2} & \ph{SRC_TF_SRDI_96_ELEC} & \ph{SRC_TF_SRDI_96_EXCH} & \ph{SRC_TF_SRDI_96_TRAF} & \ph{SRC_TF_SRDI_96_WEA} & \ph{SRC_TF_SRDI_96_MACRO} \\
    \midrule
    \rowcolor{bestgray}\introact{} & \ph{SRC_TF_OURS_96_ETTH1} & \ph{SRC_TF_OURS_96_ETTH2} & \ph{SRC_TF_OURS_96_ETTM1} & \ph{SRC_TF_OURS_96_ETTM2} & \ph{SRC_TF_OURS_96_ELEC} & \ph{SRC_TF_OURS_96_EXCH} & \ph{SRC_TF_OURS_96_TRAF} & \ph{SRC_TF_OURS_96_WEA} & \ph{SRC_TF_OURS_96_MACRO} \\
    \midrule
    Catalog Oracle & \oracle{\ph{SRC_TF_ORACLE_96_ETTH1}} & \oracle{\ph{SRC_TF_ORACLE_96_ETTH2}} & \oracle{\ph{SRC_TF_ORACLE_96_ETTM1}} & \oracle{\ph{SRC_TF_ORACLE_96_ETTM2}} & \oracle{\ph{SRC_TF_ORACLE_96_ELEC}} & \oracle{\ph{SRC_TF_ORACLE_96_EXCH}} & \oracle{\ph{SRC_TF_ORACLE_96_TRAF}} & \oracle{\ph{SRC_TF_ORACLE_96_WEA}} & \oracle{\ph{SRC_TF_ORACLE_96_MACRO}} \\
    \bottomrule
  \end{tabular}}
\end{table}

\begin{table}[h]
  \caption{Per-source MASE on TimesFM at $10\%$ severity, $H=192$. The macro column is the
  equal-weight average over the eight sources and is the quantity reported in
  Table~\ref{tab:main}. Lower is better everywhere.}
  \label{tab:app-src-tf-192}
  \centering
  \small
  \setlength{\tabcolsep}{3pt}
\resizebox{\textwidth}{!}{\begin{tabular}{lcccccccc c}
    \toprule
    Method & ETTh1 & ETTh2 & ETTm1 & ETTm2 & Electricity & Exchange & Traffic & Weather & Macro $\downarrow$ \\
    \midrule
    Native KEEP & \ph{SRC_TF_KEEP_192_ETTH1} & \ph{SRC_TF_KEEP_192_ETTH2} & \ph{SRC_TF_KEEP_192_ETTM1} & \ph{SRC_TF_KEEP_192_ETTM2} & \ph{SRC_TF_KEEP_192_ELEC} & \ph{SRC_TF_KEEP_192_EXCH} & \ph{SRC_TF_KEEP_192_TRAF} & \ph{SRC_TF_KEEP_192_WEA} & \ph{SRC_TF_KEEP_192_MACRO} \\
    Best Fixed & \ph{SRC_TF_BF_192_ETTH1} & \ph{SRC_TF_BF_192_ETTH2} & \ph{SRC_TF_BF_192_ETTM1} & \ph{SRC_TF_BF_192_ETTM2} & \ph{SRC_TF_BF_192_ELEC} & \ph{SRC_TF_BF_192_EXCH} & \ph{SRC_TF_BF_192_TRAF} & \ph{SRC_TF_BF_192_WEA} & \ph{SRC_TF_BF_192_MACRO} \\
    R2-CART & \ph{SRC_TF_R2_192_ETTH1} & \ph{SRC_TF_R2_192_ETTH2} & \ph{SRC_TF_R2_192_ETTM1} & \ph{SRC_TF_R2_192_ETTM2} & \ph{SRC_TF_R2_192_ELEC} & \ph{SRC_TF_R2_192_EXCH} & \ph{SRC_TF_R2_192_TRAF} & \ph{SRC_TF_R2_192_WEA} & \ph{SRC_TF_R2_192_MACRO} \\
    T1 & \ph{SRC_TF_T1_192_ETTH1} & \ph{SRC_TF_T1_192_ETTH2} & \ph{SRC_TF_T1_192_ETTM1} & \ph{SRC_TF_T1_192_ETTM2} & \ph{SRC_TF_T1_192_ELEC} & \ph{SRC_TF_T1_192_EXCH} & \ph{SRC_TF_T1_192_TRAF} & \ph{SRC_TF_T1_192_WEA} & \ph{SRC_TF_T1_192_MACRO} \\
    TOI & \ph{SRC_TF_TOI_192_ETTH1} & \ph{SRC_TF_TOI_192_ETTH2} & \ph{SRC_TF_TOI_192_ETTM1} & \ph{SRC_TF_TOI_192_ETTM2} & \ph{SRC_TF_TOI_192_ELEC} & \ph{SRC_TF_TOI_192_EXCH} & \ph{SRC_TF_TOI_192_TRAF} & \ph{SRC_TF_TOI_192_WEA} & \ph{SRC_TF_TOI_192_MACRO} \\
    TATO & \ph{SRC_TF_TATO_192_ETTH1} & \ph{SRC_TF_TATO_192_ETTH2} & \ph{SRC_TF_TATO_192_ETTM1} & \ph{SRC_TF_TATO_192_ETTM2} & \ph{SRC_TF_TATO_192_ELEC} & \ph{SRC_TF_TATO_192_EXCH} & \ph{SRC_TF_TATO_192_TRAF} & \ph{SRC_TF_TATO_192_WEA} & \ph{SRC_TF_TATO_192_MACRO} \\
    SRDI & \ph{SRC_TF_SRDI_192_ETTH1} & \ph{SRC_TF_SRDI_192_ETTH2} & \ph{SRC_TF_SRDI_192_ETTM1} & \ph{SRC_TF_SRDI_192_ETTM2} & \ph{SRC_TF_SRDI_192_ELEC} & \ph{SRC_TF_SRDI_192_EXCH} & \ph{SRC_TF_SRDI_192_TRAF} & \ph{SRC_TF_SRDI_192_WEA} & \ph{SRC_TF_SRDI_192_MACRO} \\
    \midrule
    \rowcolor{bestgray}\introact{} & \ph{SRC_TF_OURS_192_ETTH1} & \ph{SRC_TF_OURS_192_ETTH2} & \ph{SRC_TF_OURS_192_ETTM1} & \ph{SRC_TF_OURS_192_ETTM2} & \ph{SRC_TF_OURS_192_ELEC} & \ph{SRC_TF_OURS_192_EXCH} & \ph{SRC_TF_OURS_192_TRAF} & \ph{SRC_TF_OURS_192_WEA} & \ph{SRC_TF_OURS_192_MACRO} \\
    \midrule
    Catalog Oracle & \oracle{\ph{SRC_TF_ORACLE_192_ETTH1}} & \oracle{\ph{SRC_TF_ORACLE_192_ETTH2}} & \oracle{\ph{SRC_TF_ORACLE_192_ETTM1}} & \oracle{\ph{SRC_TF_ORACLE_192_ETTM2}} & \oracle{\ph{SRC_TF_ORACLE_192_ELEC}} & \oracle{\ph{SRC_TF_ORACLE_192_EXCH}} & \oracle{\ph{SRC_TF_ORACLE_192_TRAF}} & \oracle{\ph{SRC_TF_ORACLE_192_WEA}} & \oracle{\ph{SRC_TF_ORACLE_192_MACRO}} \\
    \bottomrule
  \end{tabular}}
\end{table}

\begin{table}[h]
  \caption{Per-source MASE on Chronos-2 at $10\%$ severity, $H=96$. The macro column is the
  equal-weight average over the eight sources and is the quantity reported in
  Table~\ref{tab:main}. Lower is better everywhere.}
  \label{tab:app-src-ch2-96}
  \centering
  \small
  \setlength{\tabcolsep}{3pt}
\resizebox{\textwidth}{!}{\begin{tabular}{lcccccccc c}
    \toprule
    Method & ETTh1 & ETTh2 & ETTm1 & ETTm2 & Electricity & Exchange & Traffic & Weather & Macro $\downarrow$ \\
    \midrule
    Native KEEP & \ph{SRC_CH2_KEEP_96_ETTH1} & \ph{SRC_CH2_KEEP_96_ETTH2} & \ph{SRC_CH2_KEEP_96_ETTM1} & \ph{SRC_CH2_KEEP_96_ETTM2} & \ph{SRC_CH2_KEEP_96_ELEC} & \ph{SRC_CH2_KEEP_96_EXCH} & \ph{SRC_CH2_KEEP_96_TRAF} & \ph{SRC_CH2_KEEP_96_WEA} & \ph{SRC_CH2_KEEP_96_MACRO} \\
    Best Fixed & \ph{SRC_CH2_BF_96_ETTH1} & \ph{SRC_CH2_BF_96_ETTH2} & \ph{SRC_CH2_BF_96_ETTM1} & \ph{SRC_CH2_BF_96_ETTM2} & \ph{SRC_CH2_BF_96_ELEC} & \ph{SRC_CH2_BF_96_EXCH} & \ph{SRC_CH2_BF_96_TRAF} & \ph{SRC_CH2_BF_96_WEA} & \ph{SRC_CH2_BF_96_MACRO} \\
    R2-CART & \ph{SRC_CH2_R2_96_ETTH1} & \ph{SRC_CH2_R2_96_ETTH2} & \ph{SRC_CH2_R2_96_ETTM1} & \ph{SRC_CH2_R2_96_ETTM2} & \ph{SRC_CH2_R2_96_ELEC} & \ph{SRC_CH2_R2_96_EXCH} & \ph{SRC_CH2_R2_96_TRAF} & \ph{SRC_CH2_R2_96_WEA} & \ph{SRC_CH2_R2_96_MACRO} \\
    T1 & \ph{SRC_CH2_T1_96_ETTH1} & \ph{SRC_CH2_T1_96_ETTH2} & \ph{SRC_CH2_T1_96_ETTM1} & \ph{SRC_CH2_T1_96_ETTM2} & \ph{SRC_CH2_T1_96_ELEC} & \ph{SRC_CH2_T1_96_EXCH} & \ph{SRC_CH2_T1_96_TRAF} & \ph{SRC_CH2_T1_96_WEA} & \ph{SRC_CH2_T1_96_MACRO} \\
    TOI & \ph{SRC_CH2_TOI_96_ETTH1} & \ph{SRC_CH2_TOI_96_ETTH2} & \ph{SRC_CH2_TOI_96_ETTM1} & \ph{SRC_CH2_TOI_96_ETTM2} & \ph{SRC_CH2_TOI_96_ELEC} & \ph{SRC_CH2_TOI_96_EXCH} & \ph{SRC_CH2_TOI_96_TRAF} & \ph{SRC_CH2_TOI_96_WEA} & \ph{SRC_CH2_TOI_96_MACRO} \\
    TATO & \ph{SRC_CH2_TATO_96_ETTH1} & \ph{SRC_CH2_TATO_96_ETTH2} & \ph{SRC_CH2_TATO_96_ETTM1} & \ph{SRC_CH2_TATO_96_ETTM2} & \ph{SRC_CH2_TATO_96_ELEC} & \ph{SRC_CH2_TATO_96_EXCH} & \ph{SRC_CH2_TATO_96_TRAF} & \ph{SRC_CH2_TATO_96_WEA} & \ph{SRC_CH2_TATO_96_MACRO} \\
    SRDI & \ph{SRC_CH2_SRDI_96_ETTH1} & \ph{SRC_CH2_SRDI_96_ETTH2} & \ph{SRC_CH2_SRDI_96_ETTM1} & \ph{SRC_CH2_SRDI_96_ETTM2} & \ph{SRC_CH2_SRDI_96_ELEC} & \ph{SRC_CH2_SRDI_96_EXCH} & \ph{SRC_CH2_SRDI_96_TRAF} & \ph{SRC_CH2_SRDI_96_WEA} & \ph{SRC_CH2_SRDI_96_MACRO} \\
    \midrule
    \rowcolor{bestgray}\introact{} & \ph{SRC_CH2_OURS_96_ETTH1} & \ph{SRC_CH2_OURS_96_ETTH2} & \ph{SRC_CH2_OURS_96_ETTM1} & \ph{SRC_CH2_OURS_96_ETTM2} & \ph{SRC_CH2_OURS_96_ELEC} & \ph{SRC_CH2_OURS_96_EXCH} & \ph{SRC_CH2_OURS_96_TRAF} & \ph{SRC_CH2_OURS_96_WEA} & \ph{SRC_CH2_OURS_96_MACRO} \\
    \midrule
    Catalog Oracle & \oracle{\ph{SRC_CH2_ORACLE_96_ETTH1}} & \oracle{\ph{SRC_CH2_ORACLE_96_ETTH2}} & \oracle{\ph{SRC_CH2_ORACLE_96_ETTM1}} & \oracle{\ph{SRC_CH2_ORACLE_96_ETTM2}} & \oracle{\ph{SRC_CH2_ORACLE_96_ELEC}} & \oracle{\ph{SRC_CH2_ORACLE_96_EXCH}} & \oracle{\ph{SRC_CH2_ORACLE_96_TRAF}} & \oracle{\ph{SRC_CH2_ORACLE_96_WEA}} & \oracle{\ph{SRC_CH2_ORACLE_96_MACRO}} \\
    \bottomrule
  \end{tabular}}
\end{table}

\begin{table}[h]
  \caption{Per-source MASE on Chronos-2 at $10\%$ severity, $H=192$. The macro column is the
  equal-weight average over the eight sources and is the quantity reported in
  Table~\ref{tab:main}. Lower is better everywhere.}
  \label{tab:app-src-ch2-192}
  \centering
  \small
  \setlength{\tabcolsep}{3pt}
\resizebox{\textwidth}{!}{\begin{tabular}{lcccccccc c}
    \toprule
    Method & ETTh1 & ETTh2 & ETTm1 & ETTm2 & Electricity & Exchange & Traffic & Weather & Macro $\downarrow$ \\
    \midrule
    Native KEEP & \ph{SRC_CH2_KEEP_192_ETTH1} & \ph{SRC_CH2_KEEP_192_ETTH2} & \ph{SRC_CH2_KEEP_192_ETTM1} & \ph{SRC_CH2_KEEP_192_ETTM2} & \ph{SRC_CH2_KEEP_192_ELEC} & \ph{SRC_CH2_KEEP_192_EXCH} & \ph{SRC_CH2_KEEP_192_TRAF} & \ph{SRC_CH2_KEEP_192_WEA} & \ph{SRC_CH2_KEEP_192_MACRO} \\
    Best Fixed & \ph{SRC_CH2_BF_192_ETTH1} & \ph{SRC_CH2_BF_192_ETTH2} & \ph{SRC_CH2_BF_192_ETTM1} & \ph{SRC_CH2_BF_192_ETTM2} & \ph{SRC_CH2_BF_192_ELEC} & \ph{SRC_CH2_BF_192_EXCH} & \ph{SRC_CH2_BF_192_TRAF} & \ph{SRC_CH2_BF_192_WEA} & \ph{SRC_CH2_BF_192_MACRO} \\
    R2-CART & \ph{SRC_CH2_R2_192_ETTH1} & \ph{SRC_CH2_R2_192_ETTH2} & \ph{SRC_CH2_R2_192_ETTM1} & \ph{SRC_CH2_R2_192_ETTM2} & \ph{SRC_CH2_R2_192_ELEC} & \ph{SRC_CH2_R2_192_EXCH} & \ph{SRC_CH2_R2_192_TRAF} & \ph{SRC_CH2_R2_192_WEA} & \ph{SRC_CH2_R2_192_MACRO} \\
    T1 & \ph{SRC_CH2_T1_192_ETTH1} & \ph{SRC_CH2_T1_192_ETTH2} & \ph{SRC_CH2_T1_192_ETTM1} & \ph{SRC_CH2_T1_192_ETTM2} & \ph{SRC_CH2_T1_192_ELEC} & \ph{SRC_CH2_T1_192_EXCH} & \ph{SRC_CH2_T1_192_TRAF} & \ph{SRC_CH2_T1_192_WEA} & \ph{SRC_CH2_T1_192_MACRO} \\
    TOI & \ph{SRC_CH2_TOI_192_ETTH1} & \ph{SRC_CH2_TOI_192_ETTH2} & \ph{SRC_CH2_TOI_192_ETTM1} & \ph{SRC_CH2_TOI_192_ETTM2} & \ph{SRC_CH2_TOI_192_ELEC} & \ph{SRC_CH2_TOI_192_EXCH} & \ph{SRC_CH2_TOI_192_TRAF} & \ph{SRC_CH2_TOI_192_WEA} & \ph{SRC_CH2_TOI_192_MACRO} \\
    TATO & \ph{SRC_CH2_TATO_192_ETTH1} & \ph{SRC_CH2_TATO_192_ETTH2} & \ph{SRC_CH2_TATO_192_ETTM1} & \ph{SRC_CH2_TATO_192_ETTM2} & \ph{SRC_CH2_TATO_192_ELEC} & \ph{SRC_CH2_TATO_192_EXCH} & \ph{SRC_CH2_TATO_192_TRAF} & \ph{SRC_CH2_TATO_192_WEA} & \ph{SRC_CH2_TATO_192_MACRO} \\
    SRDI & \ph{SRC_CH2_SRDI_192_ETTH1} & \ph{SRC_CH2_SRDI_192_ETTH2} & \ph{SRC_CH2_SRDI_192_ETTM1} & \ph{SRC_CH2_SRDI_192_ETTM2} & \ph{SRC_CH2_SRDI_192_ELEC} & \ph{SRC_CH2_SRDI_192_EXCH} & \ph{SRC_CH2_SRDI_192_TRAF} & \ph{SRC_CH2_SRDI_192_WEA} & \ph{SRC_CH2_SRDI_192_MACRO} \\
    \midrule
    \rowcolor{bestgray}\introact{} & \ph{SRC_CH2_OURS_192_ETTH1} & \ph{SRC_CH2_OURS_192_ETTH2} & \ph{SRC_CH2_OURS_192_ETTM1} & \ph{SRC_CH2_OURS_192_ETTM2} & \ph{SRC_CH2_OURS_192_ELEC} & \ph{SRC_CH2_OURS_192_EXCH} & \ph{SRC_CH2_OURS_192_TRAF} & \ph{SRC_CH2_OURS_192_WEA} & \ph{SRC_CH2_OURS_192_MACRO} \\
    \midrule
    Catalog Oracle & \oracle{\ph{SRC_CH2_ORACLE_192_ETTH1}} & \oracle{\ph{SRC_CH2_ORACLE_192_ETTH2}} & \oracle{\ph{SRC_CH2_ORACLE_192_ETTM1}} & \oracle{\ph{SRC_CH2_ORACLE_192_ETTM2}} & \oracle{\ph{SRC_CH2_ORACLE_192_ELEC}} & \oracle{\ph{SRC_CH2_ORACLE_192_EXCH}} & \oracle{\ph{SRC_CH2_ORACLE_192_TRAF}} & \oracle{\ph{SRC_CH2_ORACLE_192_WEA}} & \oracle{\ph{SRC_CH2_ORACLE_192_MACRO}} \\
    \bottomrule
  \end{tabular}}
\end{table}

\subsubsection{Action-Opportunity Strata}
\label{app:opportunity}
This subsection supports \S\ref{sec:exp-exists}. For an evaluation episode $i$ the oracle
opportunity is the improvement the best admissible action could have produced over the
reference input,
\begin{equation}
  \Delta_i \;=\; \loss\!\left(F(T_{a_0}(X_i, M_i)),\, y_i\right) - \min_{a \in \calA} \loss\!\left(F(T_{a}(X_i, M_i)),\, y_i\right),
  \label{eq:opportunity}
\end{equation}
which uses the realised target and is therefore a diagnostic rather than a deployable
quantity. Episodes are partitioned into a \emph{no-op} stratum, where the reference input is
already best up to a tolerance, a \emph{low-opportunity} stratum, and a
\emph{high-opportunity} stratum. The two boundaries are fitted on the gate split and applied
unchanged to TEST, so no evaluation episode participates in choosing them.
Table~\ref{tab:app-opp-sizes} records the realised stratum sizes,
Table~\ref{tab:app-opportunity} reports source-macro MASE within each stratum under the
identical partition for every method, and Table~\ref{tab:app-opp-sens} repeats the
stratum-level MASE under alternative boundaries so that a reader can see how much of the
effect depends on where the boundary was drawn.

\begin{table}[h]
  \caption{Action-opportunity strata on TEST. Boundaries are the gate thresholds, applied
  unchanged. Counts are of parents and not of mask variants.}
  \label{tab:app-opp-sizes}
  \centering
  \small
  \setlength{\tabcolsep}{4pt}
\resizebox{\textwidth}{!}{\begin{tabular}{lcccc}
    \toprule
    Stratum & Parents & Share of episodes & Median $\Delta_i$ & 90th pct.\ $\Delta_i$ \\
    \midrule
    No-op             & \ph{STRATUM_NOOP_N}  & \ph{STRATUM_NOOP_SHARE}  & \ph{STRATUM_NOOP_MED}  & \ph{STRATUM_NOOP_P90} \\
    Low opportunity   & \ph{STRATUM_LOW_N}   & \ph{STRATUM_LOW_SHARE}   & \ph{STRATUM_LOW_MED}   & \ph{STRATUM_LOW_P90} \\
    High opportunity  & \ph{STRATUM_HIGH_N}  & \ph{STRATUM_HIGH_SHARE}  & \ph{STRATUM_HIGH_MED}  & \ph{STRATUM_HIGH_P90} \\
    \bottomrule
  \end{tabular}}
\end{table}

\begin{table}[h]
  \caption{Stratum-level MASE under the identical partition. The no-op stratum is where an
  intervention can only lose, so a method that is better only in the high-opportunity stratum
  while degrading the no-op stratum has not solved the selection problem. Lower is better.}
  \label{tab:app-opportunity}
  \centering
  \small
  \setlength{\tabcolsep}{4pt}
\resizebox{\textwidth}{!}{\begin{tabular}{lcccc}
    \toprule
    Method & No-op MASE $\downarrow$ & Low opportunity $\downarrow$ & High opportunity $\downarrow$ & Overall $\downarrow$ \\
    \midrule
    Native KEEP        & \ph{NOOP_KEEP}    & \ph{LOWOP_KEEP}    & \ph{HIGHOP_KEEP}    & \ph{OVR_KEEP} \\
    Best Fixed         & \ph{NOOP_BF}      & \ph{LOWOP_BF}      & \ph{HIGHOP_BF}      & \ph{OVR_BF} \\
    R2-CART            & \ph{NOOP_R2}      & \ph{LOWOP_R2}      & \ph{HIGHOP_R2}      & \ph{OVR_R2} \\
    T1                 & \ph{NOOP_T1}      & \ph{LOWOP_T1}      & \ph{HIGHOP_T1}      & \ph{OVR_T1} \\
    TOI                & \ph{NOOP_TOI}     & \ph{LOWOP_TOI}     & \ph{HIGHOP_TOI}     & \ph{OVR_TOI} \\
    TATO               & \ph{NOOP_TATO}    & \ph{LOWOP_TATO}    & \ph{HIGHOP_TATO}    & \ph{OVR_TATO} \\
    SRDI               & \ph{NOOP_SRDI}    & \ph{LOWOP_SRDI}    & \ph{HIGHOP_SRDI}    & \ph{OVR_SRDI} \\
    \midrule
    \rowcolor{bestgray}\introact{} & \ph{NOOP_OURS} & \ph{LOWOP_OURS} & \ph{HIGHOP_OURS} & \ph{OVR_OURS} \\
    \bottomrule
  \end{tabular}}
\end{table}

\begin{table}[h]
  \caption{Stratum-boundary sensitivity. The no-op tolerance and the low and high boundary are
  shifted by the indicated factor of the gate threshold, and every other setting is unchanged.
  MASE is source-macro averaged within the resulting strata.}
  \label{tab:app-opp-sens}
  \centering
  \small
  \setlength{\tabcolsep}{4pt}
\resizebox{\textwidth}{!}{\begin{tabular}{lcccc}
    \toprule
    Boundary factor & No-op MASE $\downarrow$ & Low MASE $\downarrow$ & High MASE $\downarrow$ & Overall MASE $\downarrow$ \\
    \midrule
    $\times 0.5$ & \ph{OPPSENS_HALF_NOOP}  & \ph{OPPSENS_HALF_LOW}  & \ph{OPPSENS_HALF_HIGH}  & \ph{OPPSENS_HALF_OVR} \\
    $\times 1.0$ & \ph{OPPSENS_ONE_NOOP}   & \ph{OPPSENS_ONE_LOW}   & \ph{OPPSENS_ONE_HIGH}   & \ph{OPPSENS_ONE_OVR} \\
    $\times 2.0$ & \ph{OPPSENS_TWO_NOOP}   & \ph{OPPSENS_TWO_LOW}   & \ph{OPPSENS_TWO_HIGH}   & \ph{OPPSENS_TWO_OVR} \\
    \bottomrule
  \end{tabular}}
\end{table}


\subsubsection{Extended Comparison}
\label{app:extended}
Table~\ref{tab:app-extended} reports the two remaining published baselines of the same type
that are not part of the core set of Table~\ref{tab:main}. Both are variable-subset imputation
methods, so they overlap in type with TOI and SRDI, and keeping them out of the core set keeps
the main comparison balanced across the control types the paper argues about. They are excluded
from every rank in the main text. TOI-VSF~\citep{hao2025toivsf} has an official implementation
and is reported in full. GIMCC~\citep{hao2025gimcc} and VIDA~\citep{liang2025vida} have no
public implementation, and a reimplementation would not be comparable to the numbers those
papers report, so their rows are reported as unavailable rather than estimated.

If a core baseline turns out not to admit an admissible adaptation to our protocol, we report
that to the reader instead of replacing it silently, and any change to the core set is an
explicit protocol revision rather than an editorial one.

\begin{table}[h]
  \caption{Extended comparison on the main protocol, $10\%$ severity, source-macro MASE, lower
  is better. Both methods are excluded from the core average ranks of Table~\ref{tab:main}.
  A GIMCC cell is reported only if a comparable implementation is available, and is otherwise
  reported as unavailable rather than estimated.}
  \label{tab:app-extended}
  \centering
  \small
  \setlength{\tabcolsep}{4pt}
  \begin{tabular}{lccccc}
    \toprule
    Method & Bolt $\downarrow$ & TimesFM $\downarrow$ & Chronos-2 $\downarrow$ & Overall $\downarrow$ & Rank $\downarrow$ \\
    \midrule
    TOI-VSF        & \ph{TOIVSF_BOLT} & \ph{TOIVSF_TF} & \ph{TOIVSF_CH2} & \ph{TOIVSF_OVR} & \ph{TOIVSF_RANK} \\
    GIMCC          & \ph{GIMCC_BOLT}  & \ph{GIMCC_TF}  & \ph{GIMCC_CH2}  & \ph{GIMCC_OVR}  & \ph{GIMCC_RANK} \\
    \bottomrule
  \end{tabular}
\end{table}

\subsubsection{Per-Pattern Results}
\label{app:pattern}
Table~\ref{tab:app-pattern-res} separates the four patterns. Point missing is expected to
be the easiest regime and tail missing the hardest, because the missing block is adjacent
to the forecast origin and therefore removes exactly the observations a forecaster would
rely on most. Any claim about a pattern is a claim about this table, not about the pooled
main table.

\begin{table}[h]
  \caption{Per-pattern MASE at $10\%$ severity, source-macro averaged over the three
  backbones and both horizons. The last column is the average rank over the four patterns.}
  \label{tab:app-pattern-res}
  \centering
  \small
  \setlength{\tabcolsep}{4pt}
\resizebox{\textwidth}{!}{\begin{tabular}{lccccc}
    \toprule
    Method & P1 Point & P2 Target Block & P3 Shared Block & P4 Tail & Avg.\ Rank $\downarrow$ \\
    \midrule
    \textsc{Native Keep} & \ph{P1_KEEP} & \ph{P2_KEEP} & \ph{P3_KEEP} & \ph{P4_KEEP} & \ph{PR_KEEP} \\
    \textsc{Best Fixed}  & \ph{P1_BF}   & \ph{P2_BF}   & \ph{P3_BF}   & \ph{P4_BF}   & \ph{PR_BF} \\
    R2-CART              & \ph{P1_R2}   & \ph{P2_R2}   & \ph{P3_R2}   & \ph{P4_R2}   & \ph{PR_R2} \\
    T1 & \ph{P1_T1} & \ph{P2_T1} & \ph{P3_T1} & \ph{P4_T1} & \ph{PR_T1} \\
    TOI                   & \ph{P1_TOI}& \ph{P2_TOI}& \ph{P3_TOI}& \ph{P4_TOI}& \ph{PR_TOI} \\
    TATO & \ph{P1_TATO} & \ph{P2_TATO} & \ph{P3_TATO} & \ph{P4_TATO} & \ph{PR_TATO} \\
    SRDI                  & \ph{P1_SRDI}   & \ph{P2_SRDI}   & \ph{P3_SRDI}   & \ph{P4_SRDI}   & \ph{PR_SRDI} \\
    \midrule
    \rowcolor{bestgray}\introact{} & \ph{P1_OURS} & \ph{P2_OURS} & \ph{P3_OURS} & \ph{P4_OURS} & \ph{PR_OURS} \\
    \bottomrule
  \end{tabular}}
\end{table}


\section{Full Severity Results}
\label{app:severity-full}
\subsubsection{Severity Tables}
\label{app:severity}
Tables~\ref{tab:app-sev10}--\ref{tab:app-sev50} give the full severity sweep for the reduced
but type-complete set, and all three tables use the same seven methods. No method is retuned
between severities. \introact{} reuses the $k$ and $\beta$ frozen on the gate, and every
baseline reuses the models and rules fixed on TRAIN. The $10\%$ column is the main-text
condition. The $30\%$ and $50\%$ columns test whether one frozen configuration holds across
the registered severity grid, and they are not a generalisation test to unseen severity,
because the replay bank already samples from the same three levels. SRDI is the one published
method outside this reduced set, and it is compared at $10\%$ only, in Table~\ref{tab:main} and
Appendix~\ref{app:extended}.

\begin{table}[h]
  \caption{Complete results at $10\%$ severity for the reduced but type-complete set, all four metrics, source-macro averaged. Ranks are over deployable methods only and exclude the diagnostic row. This is the same seven-method set that Table~\ref{tab:robust} summarises, and the two tables are computed from the same records.}
  \label{tab:app-sev10}
  \centering
  \small
  \setlength{\tabcolsep}{4pt}
\resizebox{\textwidth}{!}{\begin{tabular}{lccccc}
    \toprule
    Method & MASE $\downarrow$ & MSE $\downarrow$ & MAE $\downarrow$ & RMSE $\downarrow$ & Rank $\downarrow$ \\
    \midrule
    \textsc{Native Keep} & \ph{SEV10-KEEP-MASE} & \ph{SEV10-KEEP-MSE} & \ph{SEV10-KEEP-MAE} & \ph{SEV10-KEEP-RMSE} & \ph{SEV10-KEEP-RANK} \\
    \textsc{Best Fixed} & \ph{SEV10-BF-MASE} & \ph{SEV10-BF-MSE} & \ph{SEV10-BF-MAE} & \ph{SEV10-BF-RMSE} & \ph{SEV10-BF-RANK} \\
    R2-CART & \ph{SEV10-R2-MASE} & \ph{SEV10-R2-MSE} & \ph{SEV10-R2-MAE} & \ph{SEV10-R2-RMSE} & \ph{SEV10-R2-RANK} \\
    T1 & \ph{SEV10-T1-MASE} & \ph{SEV10-T1-MSE} & \ph{SEV10-T1-MAE} & \ph{SEV10-T1-RMSE} & \ph{SEV10-T1-RANK} \\
    TOI & \ph{SEV10-TOI-MASE} & \ph{SEV10-TOI-MSE} & \ph{SEV10-TOI-MAE} & \ph{SEV10-TOI-RMSE} & \ph{SEV10-TOI-RANK} \\
    TATO & \ph{SEV10-TATO-MASE} & \ph{SEV10-TATO-MSE} & \ph{SEV10-TATO-MAE} & \ph{SEV10-TATO-RMSE} & \ph{SEV10-TATO-RANK} \\
    \midrule
    \rowcolor{bestgray}\introact{} & \ph{SEV10-OURS-MASE} & \ph{SEV10-OURS-MSE} & \ph{SEV10-OURS-MAE} & \ph{SEV10-OURS-RMSE} & \ph{SEV10-OURS-RANK} \\
    \midrule
    \oracle{Catalog Oracle (diagnostic)} & \oracle{\ph{SEV10-ORACLE-MASE}} & \oracle{\ph{SEV10-ORACLE-MSE}} & \oracle{\ph{SEV10-ORACLE-MAE}} & \oracle{\ph{SEV10-ORACLE-RMSE}} & \oracle{---} \\
    \bottomrule
  \end{tabular}}
\end{table}

\begin{table}[h]
  \caption{Complete results at $30\%$ severity for the same seven-method set, all four metrics, source-macro averaged. The configuration is identical to the $10\%$ table and nothing is retuned.}
  \label{tab:app-sev30}
  \centering
  \small
  \setlength{\tabcolsep}{4pt}
\resizebox{\textwidth}{!}{\begin{tabular}{lccccc}
    \toprule
    Method & MASE $\downarrow$ & MSE $\downarrow$ & MAE $\downarrow$ & RMSE $\downarrow$ & Rank $\downarrow$ \\
    \midrule
    \textsc{Native Keep} & \ph{SEV30-KEEP-MASE} & \ph{SEV30-KEEP-MSE} & \ph{SEV30-KEEP-MAE} & \ph{SEV30-KEEP-RMSE} & \ph{SEV30-KEEP-RANK} \\
    \textsc{Best Fixed} & \ph{SEV30-BF-MASE} & \ph{SEV30-BF-MSE} & \ph{SEV30-BF-MAE} & \ph{SEV30-BF-RMSE} & \ph{SEV30-BF-RANK} \\
    R2-CART & \ph{SEV30-R2-MASE} & \ph{SEV30-R2-MSE} & \ph{SEV30-R2-MAE} & \ph{SEV30-R2-RMSE} & \ph{SEV30-R2-RANK} \\
    T1 & \ph{SEV30-T1-MASE} & \ph{SEV30-T1-MSE} & \ph{SEV30-T1-MAE} & \ph{SEV30-T1-RMSE} & \ph{SEV30-T1-RANK} \\
    TOI & \ph{SEV30-TOI-MASE} & \ph{SEV30-TOI-MSE} & \ph{SEV30-TOI-MAE} & \ph{SEV30-TOI-RMSE} & \ph{SEV30-TOI-RANK} \\
    TATO & \ph{SEV30-TATO-MASE} & \ph{SEV30-TATO-MSE} & \ph{SEV30-TATO-MAE} & \ph{SEV30-TATO-RMSE} & \ph{SEV30-TATO-RANK} \\
    \midrule
    \rowcolor{bestgray}\introact{} & \ph{SEV30-OURS-MASE} & \ph{SEV30-OURS-MSE} & \ph{SEV30-OURS-MAE} & \ph{SEV30-OURS-RMSE} & \ph{SEV30-OURS-RANK} \\
    \bottomrule
  \end{tabular}}
\end{table}

\begin{table}[h]
  \caption{Complete results at $50\%$ severity for the same seven-method set. This level is inside the registered grid that the replay bank samples from, so the rows show within-grid robustness of one frozen configuration and not behaviour on unseen severity.}
  \label{tab:app-sev50}
  \centering
  \small
  \setlength{\tabcolsep}{4pt}
\resizebox{\textwidth}{!}{\begin{tabular}{lccccc}
    \toprule
    Method & MASE $\downarrow$ & MSE $\downarrow$ & MAE $\downarrow$ & RMSE $\downarrow$ & Rank $\downarrow$ \\
    \midrule
    \textsc{Native Keep} & \ph{SEV50-KEEP-MASE} & \ph{SEV50-KEEP-MSE} & \ph{SEV50-KEEP-MAE} & \ph{SEV50-KEEP-RMSE} & \ph{SEV50-KEEP-RANK} \\
    \textsc{Best Fixed} & \ph{SEV50-BF-MASE} & \ph{SEV50-BF-MSE} & \ph{SEV50-BF-MAE} & \ph{SEV50-BF-RMSE} & \ph{SEV50-BF-RANK} \\
    R2-CART & \ph{SEV50-R2-MASE} & \ph{SEV50-R2-MSE} & \ph{SEV50-R2-MAE} & \ph{SEV50-R2-RMSE} & \ph{SEV50-R2-RANK} \\
    T1 & \ph{SEV50-T1-MASE} & \ph{SEV50-T1-MSE} & \ph{SEV50-T1-MAE} & \ph{SEV50-T1-RMSE} & \ph{SEV50-T1-RANK} \\
    TOI & \ph{SEV50-TOI-MASE} & \ph{SEV50-TOI-MSE} & \ph{SEV50-TOI-MAE} & \ph{SEV50-TOI-RMSE} & \ph{SEV50-TOI-RANK} \\
    TATO & \ph{SEV50-TATO-MASE} & \ph{SEV50-TATO-MSE} & \ph{SEV50-TATO-MAE} & \ph{SEV50-TATO-RMSE} & \ph{SEV50-TATO-RANK} \\
    \midrule
    \rowcolor{bestgray}\introact{} & \ph{SEV50-OURS-MASE} & \ph{SEV50-OURS-MSE} & \ph{SEV50-OURS-MAE} & \ph{SEV50-OURS-RMSE} & \ph{SEV50-OURS-RANK} \\
    \bottomrule
  \end{tabular}}
\end{table}

\subsubsection{Robustness Trend Figure}
\label{app:robust-fig}
Figure~\ref{fig:robust} shows the severity trend summarised by Table~\ref{tab:robust}. Each
panel is a different frozen backbone. The axes are identical across panels so that the
panels can be compared directly.

\begin{figure}[h]
  \centering
  \includegraphics[width=0.86\textwidth]{figure/fig5_robustness_missingness.pdf}
  \caption{Relative MASE improvement over \textsc{Native Keep} as missingness severity
  increases, one line per method, one panel per frozen backbone. All methods reuse the
  models and rules frozen on TRAIN, with no per-severity retuning.}
  \label{fig:robust}
\end{figure}



\section{Reconstruction and Utility Diagnostics}
\label{app:diagnostics}
\subsubsection{Reconstruction versus Forecasting Utility}
\label{app:recutils}
This subsection supports \S\ref{sec:exp-exists}. Reconstruction error is only defined for
methods that emit an explicit estimate of the hidden positions, so \textsc{Keep}, a
transformation pipeline such as TATO, the selected intervention of \introact{}, and the
external reference ChannelTokenFormer are not placed on a common reconstruction scale. The
table is therefore restricted to the core methods that do emit reconstructed values, and the
ranking analysis is performed within each parent rather than on a pooled sample, because
parents differ in scale and in difficulty. A method that does not reconstruct is not ranked
against a method that does, and \introact{} never enters this ranking because it returns a
forecast from an executed input rather than a reconstruction. Table~\ref{tab:app-recutils}
gives the per-method numbers. Reconstruction error uses the artificially hidden positions,
which are known because the mask is synthetic, and forecast gain is measured against
\textsc{Native Keep} on the same backbone. The discordant rate is the fraction of method pairs
whose within-parent reconstruction and forecasting rankings disagree.

A \emph{Reconstruction Oracle} is included as a diagnostic. It selects the action with the
smallest reconstruction error and therefore reads the hidden values. It is not deployable, it
is excluded from every rank, and it is reported only to show how much of the achievable
forecasting improvement a reconstruction criterion would capture if it were applied with full
knowledge of the missing entries. \ph{RECON_ORACLE_MASE}.

\begin{table}[h]
  \caption{Paired reconstruction-versus-utility table, restricted to the core methods that
  emit an explicit reconstruction of the hidden positions. The Spearman $\rho$ row compares the
  reconstruction ranking with the forecasting ranking within each backbone. A $\rho$ close to
  zero means the two rankings carry different information.}
  \label{tab:app-recutils}
  \centering
  \small
  \setlength{\tabcolsep}{4pt}
\resizebox{\textwidth}{!}{\begin{tabular}{lccccc}
    \toprule
    Method & Rec.\ MSE $\downarrow$ & Rec.\ MAE $\downarrow$ & Forecast MASE $\downarrow$
           & Forecast Gain $\uparrow$ & Discordant Rate $\downarrow$ \\
    \midrule
    T1                 & \ph{REC_T1_MSE}    & \ph{REC_T1_MAE}    & \ph{REC_T1_MASE}    & \ph{REC_T1_GAIN}    & \multirow{3}{*}{\ph{DISCORDANT_RATE}} \\
    TOI                & \ph{REC_TOI_MSE}   & \ph{REC_TOI_MAE}   & \ph{REC_TOI_MASE}   & \ph{REC_TOI_GAIN}   & \\
    SRDI               & \ph{REC_SRDI_MSE}  & \ph{REC_SRDI_MAE}  & \ph{REC_SRDI_MASE}  & \ph{REC_SRDI_GAIN}  & \\
    \midrule
    \multicolumn{6}{l}{\footnotesize Spearman $\rho$ between reconstruction rank and
    forecasting rank: Bolt \ph{RHO_BOLT}, TimesFM \ph{RHO_TF}, Chronos-2 \ph{RHO_CH2}} \\
    \bottomrule
  \end{tabular}}
\end{table}

\subsubsection{Governance Diagnostics}
\label{app:governance-fig}
Figure~\ref{fig:governance} visualises the diagnostics tabulated in
Table~\ref{tab:harm}. Panel~(b) is a sign-agreement check rather than a calibration
curve: the ranking score is used only to order actions, and we do not claim that its magnitude
is a calibrated probability.

\begin{figure}[h]
  \centering
  \includegraphics[width=0.84\textwidth]{figure/fig4_governance_diagnostics.pdf}
  \caption{\textbf{(a)}~Harmful loss by method. \textbf{(b)}~Realised utility against the
  binned ranking score with 95\% intervals. Panel~(b) is a sign agreement check, not a
  calibration claim.}
  \label{fig:governance}
\end{figure}



\section{Ablations and Stability}
\label{app:ablations-full}
\subsection{Ablation Details}
\label{app:ablation-details}
This section supports \S\ref{sec:exp-ablation}. Every variant reads the same catalog
prediction cache, so the rows differ only in how the stored utilities are summarised and
thresholded. Table~\ref{tab:app-ablation} reports the full outcome vector behind the compact
Table~\ref{tab:ablation}. The variant definitions follow \S\ref{sec:exp-ablation} exactly.

\begin{table}[h]
  \caption{Ablation, full outcome vector. The Reconstruction Oracle reads the hidden values,
  is a diagnostic rather than a deployable row, and is excluded from every rank. Calls counts
  forecasting-backbone calls per request. Harmful loss is measured relative to
  \textsc{Keep}.}
  \label{tab:app-ablation}
  \centering
  \small
  \setlength{\tabcolsep}{4pt}
\resizebox{\textwidth}{!}{\begin{tabular}{lcccccccc}
    \toprule
    Variant & MASE $\downarrow$ & MAE $\downarrow$ & RMSE $\downarrow$ & MSE $\downarrow$ & Intervention rate & Conditional HIR $\downarrow$ & Harmful loss $\downarrow$ & Calls $\downarrow$ \\
    \midrule
    Full \introact{}            & \ph{AF_FULL_MASE} & \ph{AF_FULL_MAE} & \ph{AF_FULL_RMSE} & \ph{AF_FULL_MSE} & \ph{AF_FULL_IR} & \ph{AF_FULL_CHIR} & \ph{AF_FULL_HL} & \ph{AF_FULL_CALLS} \\
    A1 Global utility           & \ph{AF_A1_MASE}   & \ph{AF_A1_MAE}   & \ph{AF_A1_RMSE}   & \ph{AF_A1_MSE}   & \ph{AF_A1_IR}   & \ph{AF_A1_CHIR}   & \ph{AF_A1_HL}   & \ph{AF_A1_CALLS} \\
    A2 w/o intervention state   & \ph{AF_A2_MASE}   & \ph{AF_A2_MAE}   & \ph{AF_A2_RMSE}   & \ph{AF_A2_MSE}   & \ph{AF_A2_IR}   & \ph{AF_A2_CHIR}   & \ph{AF_A2_HL}   & \ph{AF_A2_CALLS} \\
    A3 w/o reference-forecast state & \ph{AF_A3_MASE} & \ph{AF_A3_MAE} & \ph{AF_A3_RMSE} & \ph{AF_A3_MSE}   & \ph{AF_A3_IR}   & \ph{AF_A3_CHIR}   & \ph{AF_A3_HL}   & \ph{AF_A3_CALLS} \\
    A4 w/o conservative gate    & \ph{AF_A4_MASE}   & \ph{AF_A4_MAE}   & \ph{AF_A4_RMSE}   & \ph{AF_A4_MSE}   & \ph{AF_A4_IR}   & \ph{AF_A4_CHIR}   & \ph{AF_A4_HL}   & \ph{AF_A4_CALLS} \\
    A5 Parametric utility predictor & \ph{AF_A5_MASE} & \ph{AF_A5_MAE} & \ph{AF_A5_RMSE} & \ph{AF_A5_MSE} & \ph{AF_A5_IR} & \ph{AF_A5_CHIR} & \ph{AF_A5_HL} & \ph{AF_A5_CALLS} \\
    \midrule
    Reconstruction Oracle (diagnostic) & \oracle{\ph{AF_ORACLE_MASE}} & \oracle{\ph{AF_ORACLE_MAE}} & \oracle{\ph{AF_ORACLE_RMSE}} & \oracle{\ph{AF_ORACLE_MSE}} & \oracle{\ph{AF_ORACLE_IR}} & \oracle{\ph{AF_ORACLE_CHIR}} & \oracle{\ph{AF_ORACLE_HL}} & \oracle{--} \\
    \bottomrule
  \end{tabular}}
\end{table}

\subsection{Mask-Realisation Stability}
\label{app:seeds}
The main experiment uses one deterministic mask per parent. This appendix repeats the
comparison under two further deterministic mask realisations of the same parents, so that a
reader can see whether a conclusion depends on one particular realisation of the deletion
pattern. The three realisations are fixed before the run and are never treated as three
independent samples of parents: the statistical unit stays the parent, and the three
realisations are read as a stability check on the same parent set. Significance testing is
never performed across them.

The subset is fixed before the run: all eight sources, patterns P1 to P4, $10\%$ severity,
$H=96$, on Bolt and Chronos-2, and the comparison is restricted to \textsc{Keep}, R2-CART,
and \introact{}. Reporting the spread over three realisations is a stability diagnostic and
is not an additional estimate of variance.

\begin{table}[h]
  \caption{Mask-realisation stability. Each row is one deterministic mask realisation on the
  pre-registered subset. The spread across realisations is a stability diagnostic and is not
  an additional sample of parents.}
  \label{tab:app-seeds}
  \centering
  \small
  \setlength{\tabcolsep}{4pt}
\resizebox{\textwidth}{!}{\begin{tabular}{llccc}
    \toprule
    Backbone & Mask realisation & MASE $\downarrow$ & $\Delta$ vs \textsc{Keep} $\uparrow$ & Intervention rate \\
    \midrule
    Bolt      & primary seed        & \ph{SEED1_MASE} & \ph{SEED1_DELTA} & \ph{SEED1_IR} \\
    Bolt      & second realisation  & \ph{SEED2_MASE} & \ph{SEED2_DELTA} & \ph{SEED2_IR} \\
    Bolt      & third realisation   & \ph{SEED3_MASE} & \ph{SEED3_DELTA} & \ph{SEED3_IR} \\
    \midrule
    Chronos-2 & primary seed        & \ph{SEED4_MASE} & \ph{SEED4_DELTA} & \ph{SEED4_IR} \\
    Chronos-2 & second realisation  & \ph{SEED5_MASE} & \ph{SEED5_DELTA} & \ph{SEED5_IR} \\
    Chronos-2 & third realisation   & \ph{SEED6_MASE} & \ph{SEED6_DELTA} & \ph{SEED6_IR} \\
    \midrule
    \multicolumn{2}{l}{Spread over the three realisations (Bolt)}      & \ph{SEED_SPREAD_BOLT} & --- & \ph{SEED_IR_SPREAD_BOLT} \\
    \multicolumn{2}{l}{Spread over the three realisations (Chronos-2)} & \ph{SEED_SPREAD_CH2}  & --- & \ph{SEED_IR_SPREAD_CH2} \\
    \bottomrule
  \end{tabular}}
\end{table}

\section{Cost, Failures, and Memory}
\label{app:cost-full}
\subsection{Runtime and Failure Audit}
\label{app:calls}
Table~\ref{tab:app-efficiency} gives the full cost breakdown behind the summary in
\S\ref{sec:exp-ablation}. Offline cost and online cost are kept in separate columns because
they are not directly comparable: the offline replay bank is built once per backbone revision,
per source, and per horizon, whereas the online columns are per request. Within each column a
lower value is better, and the two columns are read separately rather than combined into a
single score. Latency percentiles are measured per request and exclude cold start, which is
reported on its own line. Candidate materialisation is a separate line because it is paid
before the decision is made and is not a forecasting-backbone call.

\begin{table}[t]
  \caption{Deployment cost for the methods that carry a claim, offline and online kept
  separate. Latency percentiles are per request and exclude cold start, reported on its own
  line. Backbone calls counts forecasting-backbone calls per request and excludes candidate
  materialisation, which is reported in its own column. The two extended methods of
  Appendix~\ref{app:extended} are not given separate cost rows because they carry no claim.}
  \label{tab:app-efficiency}
  \centering
  \small
  \setlength{\tabcolsep}{4pt}
\resizebox{\textwidth}{!}{\begin{tabular}{lcccccc}
    \toprule
    Method & Offline train/search & Candidate constr. & Backbone calls/req.\ & Mean lat.\ $\downarrow$ & P95 $\downarrow$ & Max $\downarrow$ \\
    \midrule
    T1                 & \ph{T1_OFF}     & \ph{T1_CAND}     & \ph{T1_CALLS}     & \ph{T1_MEAN}     & \ph{T1_P95}     & \ph{T1_MAX} \\
    TOI                & \ph{TOI_OFF}    & \ph{TOI_CAND}    & \ph{TOI_CALLS}    & \ph{TOI_MEAN}    & \ph{TOI_P95}    & \ph{TOI_MAX} \\
    TATO               & \ph{TATO_OFF}   & \ph{TATO_CAND}   & \ph{TATO_CALLS}   & \ph{TATO_MEAN}   & \ph{TATO_P95}   & \ph{TATO_MAX} \\
    SRDI               & \ph{SRDI_OFF}   & \ph{SRDI_CAND}   & \ph{SRDI_CALLS}   & \ph{SRDI_MEAN}   & \ph{SRDI_P95}   & \ph{SRDI_MAX} \\
    ChannelTokenFormer & \ph{CTF_OFF}    & \ph{CTF_CAND}    & \ph{CTF_CALLS}    & \ph{CTF_MEAN}    & \ph{CTF_P95}    & \ph{CTF_MAX} \\
    \rowcolor{bestgray}\introact{} & \ph{OURS_OFF} & \ph{OURS_CAND} & \ph{OURS_CALLS} & \ph{OURS_MEAN} & \ph{OURS_P95} & \ph{OURS_MAX} \\
    \midrule
    \multicolumn{7}{l}{\footnotesize Cold start: \ph{COLDSTART} \quad
    Failure/timeout share: \ph{FAILRATE}} \\
    \bottomrule
  \end{tabular}}
\end{table}

Table~\ref{tab:app-calls} separates offline from online cost and separates successes from
failures. Failures and timeouts are never counted as successes and are never dropped from the
denominator. A request that fails to produce a forecast is reported as a failure rather than
as a zero-cost success. A request costs one forecasting-backbone call for the reference action
and, when an intervention is executed, one further call for the selected action, so the online
count is between one and two backbone calls. When the decision returns \textsc{Keep} the input
is left untouched and the reference call is the returned forecast. Candidate materialisation
and scoring are booked on their own line because they are paid before the decision and do not
touch the backbone.

\begin{table}[h]
  \caption{Backbone call accounting. Latency is reported per request and excludes cold start,
  which is reported on its own line. The offline replay-bank build and the one-off selector fit
  are amortised and reported separately. Candidate materialisation and scoring are booked on
  their own line and are not forecasting-backbone calls.}
  \label{tab:app-calls}
  \centering
  \small
  \setlength{\tabcolsep}{4pt}
\resizebox{\textwidth}{!}{\begin{tabular}{lccc}
    \toprule
    Outcome & Requests & Share & Backbone calls per request \\
    \midrule
    \textsc{Keep} returned (input untouched)   & \ph{CALLS_KEEP_N}    & \ph{CALLS_KEEP_SHARE}    & 1 \\
    Intervention executed (input changed)      & \ph{CALLS_ACT_N}     & \ph{CALLS_ACT_SHARE}     & 2 \\
    Failed before any forecast                 & \ph{CALLS_FAIL_N}    & \ph{CALLS_FAIL_SHARE}    & \ph{CALLS_FAIL_COST} \\
    Timed out after the forecast               & \ph{CALLS_TIMEOUT_N} & \ph{CALLS_TIMEOUT_SHARE} & \ph{CALLS_TIMEOUT_COST} \\
    \midrule
    Candidate materialisation and scoring (per request) & \ph{CALLS_CAND_N} & --- & 0 \\
    Offline replay-bank build (amortised)      & --- & --- & \ph{CALLS_OFFLINE} \\
    Offline selector fit (amortised)           & --- & --- & \ph{CALLS_SELECTOR} \\
    Cold start (first request of a process)    & \ph{CALLS_COLD_N} & --- & \ph{CALLS_COLD_LATENCY} \\
    \bottomrule
  \end{tabular}}
\end{table}

\section{Reproducibility and Artefact Schema}
\label{app:repro-full}
\subsubsection{Reproducibility Details}
\label{app:reprodetails}
Each record in the result files contains: source, parent, origin, horizon, pattern,
severity, mask hash, backbone, backbone revision, method, action, input hash, prediction
hash, forecast, target reference id, MASE, MSE, MAE, RMSE, reconstruction MSE where
applicable, runtime components, failure, fallback, code SHA, and the resolved configuration.
Results are written to the logical result groups listed in Table~\ref{tab:app-groups}, each in
CSV, JSON, and Markdown form, with raw predictions retained alongside the aggregates. The
group names describe the stage that produced them and do not encode a version number, so that
a later revision of the protocol cannot be confused with the present one.

\begin{table}[h]
  \caption{Result groups written by the pipeline. Every number in the main text and in the
  appendices is generated from one of these groups and nothing is transcribed by hand.}
  \label{tab:app-groups}
  \centering
  \small
  \setlength{\tabcolsep}{4pt}
\resizebox{\textwidth}{!}{\begin{tabular}{lll}
    \toprule
    Result group & Contents & Consumed by \\
    \midrule
    \texttt{protocol}       & frozen protocol, source manifests, split boundaries, TEST manifest & Appendix~\ref{app:splits} \\
    \texttt{replay\_bank}   & replay records $(z_{i,a}, a, g_{i,a})$ with input and prediction hashes & \S\ref{sec:method-replay} \\
    \texttt{baselines}      & baseline training runs, resolved hyperparameters, checkpoints & \S\ref{sec:exp-main} \\
    \texttt{train\_eval}    & TRAIN-eval acceptance check before freezing & Appendix~\ref{app:splits} \\
    \texttt{main\_results}  & main matrix rows, all methods, all cells & Tables~\ref{tab:main}, \ref{tab:app-src-bolt-96}--\ref{tab:app-src-ch2-192} \\
    \texttt{ablations}      & the six ablation variants plus the Reconstruction Oracle diagnostic & Table~\ref{tab:ablation}, \ref{tab:app-ablation} \\
    \texttt{robustness}     & severity, pattern, and mask-realisation sweeps & Tables~\ref{tab:robust}, \ref{tab:app-sev10}--\ref{tab:app-sev50}, \ref{tab:app-seeds} \\
    \texttt{diagnostics}    & reconstruction and utility diagnostics, rank agreement & Appendix~\ref{app:diagnostics} \\
    \texttt{cost\_audit}    & call counts, latency, failures, timeouts, memory & Table~\ref{tab:app-efficiency}, Appendix~\ref{app:calls} \\
    \bottomrule
  \end{tabular}}
\end{table}

The protocol is frozen before the main experiment and identified by commit hash. Any change to
the action catalog, the missingness protocol, the mask derivation rule, the split boundaries,
the state feature list, or the failure policy invalidates the freeze and requires a new
configuration identifier. Calibration and TEST labels are not read during development, and no
result in this paper is selected on them.

\paragraph{Software and hardware.} Forecasts are produced with the released checkpoints of
Chronos-Bolt, TimesFM, and Chronos-2 at the revisions listed in Table~\ref{tab:app-repro}.
The surrounding code is Python \ph{PYVERSION} with NumPy \ph{NUMPYVERSION}, pandas
\ph{PANDASVERSION}, scikit-learn \ph{SKLEARNVERSION}, and PyTorch \ph{TORCHVERSION}. Runs are
executed on \ph{HARDWARE} in \ph{PRECISION}. The development-cycle numbers of
Appendix~\ref{app:development} were produced on a different machine and are not mixed with the
present results.

\begin{table}[h]
  \caption{Reproducibility anchors. Backbone revisions are the exact released checkpoints used
  for every forecast in this paper. Seeds are the deterministic seeds that fix the mask
  derivation, the replay-record ordering, and the selector fit.}
  \label{tab:app-repro}
  \centering
  \small
  \setlength{\tabcolsep}{4pt}
  \begin{tabular}{@{}l p{0.70\textwidth}@{}}
    \toprule
    Item & Value \\
    \midrule
    Chronos-Bolt revision   & \ph{REV_BOLT} \\
    TimesFM revision        & \ph{REV_TF} \\
    Chronos-2 revision      & \ph{REV_CH2} \\
    Python and libraries    & \ph{PYVERSION}, NumPy \ph{NUMPYVERSION}, pandas \ph{PANDASVERSION}, scikit-learn \ph{SKLEARNVERSION}, PyTorch \ph{TORCHVERSION} \\
    Hardware and precision  & \ph{HARDWARE}, \ph{PRECISION} \\
    Mask derivation seed    & \ph{SEED_MASK} \\
    Replay-record order seed & \ph{SEED_ORDER} \\
    Selector fit seed       & \ph{SEED_SELECTOR} \\
    Protocol commit         & \ph{PROTOCOL_COMMIT} \\
    TEST manifest           & \ph{TEST_MANIFEST} \\
    \bottomrule
  \end{tabular}
\end{table}

The TEST manifest is the authoritative record of which parents and origins were held out. It
is written once, before any result is computed, and every table in this paper is generated
against that manifest rather than against a split reconstructed after the fact.

\section*{Supplementary Exploratory Material}
\label{app:supplementary}
\subsection{Development Negative Results}
\label{app:development}
This section records what was tried earlier in this line of work and rejected. It is included
because the final formulation is easier to read against the alternatives it replaced, and
because omitting it would misrepresent how the method was reached. Nothing here is independent
confirmation of any method, and no number here is comparable with any table in the main text.

An earlier configuration tried to improve the action ranking by projecting estimated utility
vectors onto a convex feasible set before thresholding them. The projection reduced the squared
error of the estimated vector while failing to improve the final forecasting metric, and on one
backbone it degraded it. That outcome is the reason the present method ranks actions with a
local mean and a dispersion penalty and then applies a single threshold, rather than
calibrating a continuous score and reading a decision off it.

Two consequences of that cycle carry into the present paper. The strongest simple selector
available at the time, which is R2-CART in the present terminology, was competitive with the
more elaborate candidate, and on one backbone better than it, which is why R2-CART is carried
forward as a mandatory control in every main-text table. Reducing estimation error on a
continuous score was also not sufficient to improve discrete selection, which is why the
conservative gate is evaluated directly on harmful intervention rate and harmful loss rather
than on MASE alone.

The measurements behind that decision come from a different protocol, with a smaller
development split, fewer sources, and different budget accounting, so they are not reproduced
here as numbers. They are recorded in the released artefact ledger and must not be compared
numerically with any table in this paper.

\end{document}
